\documentclass[11pt]{extarticle}
\usepackage{amsmath,amssymb,amsfonts,amsthm}
\usepackage{mathtools}
\usepackage{enumitem}
\usepackage[margin=1in]{geometry}
\usepackage{hyperref}
\hypersetup{hidelinks}

\newtheorem{theorem}{Theorem}[section]
\newtheorem{lemma}[theorem]{Lemma}
\newtheorem{corollary}[theorem]{Corollary}
\newtheorem{conjecture}[theorem]{Conjecture}
\theoremstyle{definition}
\newtheorem{definition}[theorem]{Definition}
\newtheorem{proposition}[theorem]{Proposition}
\newtheorem{example}[theorem]{Example}
\theoremstyle{remark}
\newtheorem{remark}[theorem]{Remark}

\numberwithin{equation}{section}

%------------------------------------------------
% Notation
%------------------------------------------------
\newcommand{\F}{\mathbb F}
\newcommand{\N}{\mathbb N}
\newcommand{\Z}{\mathbb Z}
\newcommand{\Oq}{O_q}
\newcommand{\Aq}{\mathcal A_q}
\newcommand{\Aqfib}[1]{\mathcal A_{q,#1}}
\newcommand{\BBqquot}[2]{\widetilde{\mathcal A}_{q,#1}^{\{#2\}}}
\newcommand{\Uplus}{U_q^+(\widehat{\mathfrak{sl}}_2)}
\newcommand{\Vsh}{\mathbb V}
\newcommand{\gr}{\operatorname{gr}}
\newcommand{\init}{\operatorname{in}}
\newcommand{\sh}{\mathrm{sh}}
\newcommand{\cW}{\mathcal W}
\newcommand{\cG}{\mathcal G}
\newcommand{\ctG}{\widetilde{\mathcal G}}
\newcommand{\XiC}{\Xi}
\newcommand{\Gt}{\widetilde{\cG}}
\newcommand{\Zvee}{\mathcal Z^{\vee}}
\newcommand{\Zconv}{\mathcal Z}
\newcommand{\qcomm}[2]{[ #1,#2 ]_q}
\newcommand{\comm}[2]{[ #1,#2 ]}
\DeclareMathOperator{\Cent}{Cent}

\begin{document}

\title{PBW bases and centralisers for the $q$-Onsager algebra}

\author{Haoran Zhu\thanks{Division of Mathematical Sciences,
Nanyang Technological University, 21 Nanyang Link, Singapore 637371.
Email: \texttt{zhuh0031@e.ntu.edu.sg}}}

\date{}

\maketitle

\begin{abstract}
We settle five conjectures concerning PBW bases and centralisers for the
$q$-Onsager algebra.  We first prove that the Baseilhac--Kolb root vectors
give a PBW basis in every linear order whenever $q$ is not a root of unity,
removing the previous transcendence hypothesis.  We next establish twelve
PBW bases in the alternating generators and show that they persist under
arbitrary scalar central specialisation of the alternating central
extension.  This proves conjectures of Terwilliger and of Baseilhac and Belliard.  Finally, we
determine the centralisers of the negative and imaginary alternating
subalgebras and show that all four single-family alternating polynomial
subalgebras are maximal commutative.  The proofs combine explicit
straightening for Damiani root vectors, the degeneration
$\gr\Oq\simeq\Uplus$, and large-index triangular crossing arguments.
\end{abstract}

{\noindent\itshape Keywords:} $q$-Onsager algebra, PBW basis, alternating
generators, root vectors, central specialisation, centraliser.

\medskip

{\noindent\itshape Mathematics Subject Classification:}
17B37, 17B67, 81R50.

\section{Introduction}
\label{sec:introduction}

The Onsager algebra is one of the classical algebraic structures behind
exactly solvable lattice models.  It entered mathematical physics through
Onsager's solution of the two-dimensional Ising model
\cite{Onsager1944}, while the Dolan--Grady relations
\cite{DolanGrady1982} gave a concise algebraic mechanism for producing the
commuting quantities that underlie Onsager-type integrability.  The
$q$-Onsager algebra $\Oq$ is a $q$-deformation of this structure.  It appears
naturally in boundary quantum integrable systems
\cite{Sklyanin1988,Baseilhac2005,BaseilhacKoizumi2005,
BaseilhacBelliard2010,BaseilhacBelliard2013}, in the theory of tridiagonal
pairs and tridiagonal algebras \cite{Terwilliger2001}, and as a fundamental
example of a quantum symmetric-pair coideal subalgebra
\cite{Letzter2002,Kolb2014,LuWang2021,LuRuanWang2023}.

Two families of PBW constructions have played a central role in the study of
$\Oq$.  Baseilhac and Kolb used Lusztig-type automorphisms to construct real
and imaginary root vectors and obtained a root-vector PBW theorem
\cite{BaseilhacKolb2020}.  A second construction starts from the current
algebra introduced in the reflection-equation setting by Baseilhac and
Koizumi and by Baseilhac and Shigechi
\cite{BaseilhacKoizumi2005,BaseilhacShigechi2010}.  Terwilliger proved that
this current algebra is the alternating central extension $\Aq$ of $\Oq$
and established a standard tensor factorisation
\[
        \Aq\simeq\Oq\otimes\F[z_1,z_2,\ldots].
\]
The four generating families of $\Aq$ descend under central reduction to the
alternating generators
\[
 \{W_{-n}\}_{n\geq0},\qquad
 \{W_{n+1}\}_{n\geq0},\qquad
 \{G_{n+1}\}_{n\geq0},\qquad
 \{\widetilde G_{n+1}\}_{n\geq0}
\]
of $\Oq$ \cite{Terwilliger2021ACE,Terwilliger2022OqACE}.

The present paper addresses five conjectures arising from these two PBW
pictures.  Four of them appear in the list of future directions in
\cite[Section~16]{Terwilliger2022OqACE}.

\begin{enumerate}[label=\textup{(\roman*)},leftmargin=11mm]
\item Conjecture~16.1 asks whether the Baseilhac--Kolb root-vector PBW theorem
remains valid when the hypothesis that $q$ is transcendental is replaced by
the assumption that $q$ is not a root of unity.

\item Conjecture~16.2 predicts six PBW bases for the three alternating
families $W^-$, $\widetilde G$, and $W^+$.

\item Conjecture~16.7 predicts that the commutative algebra generated by
$W_0,W_{-1},W_{-2},\ldots$ is self-centralising in $\Oq$.

\item Conjecture~16.8 makes the analogous prediction for the algebra
generated by $\widetilde G_1,\widetilde G_2,\ldots$.

\item Baseilhac and Belliard conjectured that every central quotient of the
current algebra determined by their central values $\Delta_n=2\delta_n$
has a PBW basis in the block order $W^-<G<W^+$
\cite[Conjecture~1]{BaseilhacBelliard2017}.
\end{enumerate}

These questions are related but are not formal consequences of one another.
The root-vector problem concerns specialisation of a generic PBW theorem;
the alternating problem concerns the highest filtered components of a very
different system of generators; and the centraliser problems require a
no-cancellation argument beyond a dimension count.  Their common feature is
triangularity.  The main theme of this paper is that global degeneration and
local crossings are two manifestations of the same triangular principle.

The order of a PBW basis is relevant in applications.  Factorised formulae
for universal $R$-matrices are built from specific convex orders
\cite{Beck1994,KhoroshkinTolstoy1991,Damiani1998}; after evaluation, this dependence
is visible in LDU decompositions of $L$-operators
\cite[eqs.~(4.8) and~(4.9)]{DingFrenkel1993}.  In the coideal setting,
universal $K$-matrices and fused $K$-operators provide the corresponding
structures \cite{AppelVlaar2022,LemartheBaseilhacGainutdinov2026}.  The
alternating order conjectured by Baseilhac and Belliard has already been
used in universal TT- and TQ-relations
\cite{BaseilhacGainutdinovLemarthe2026}.  Rigorous PBW results for the
specific alternating orders are therefore needed both for the internal
structure of $\Oq$ and for order-sensitive constructions in boundary
integrability.

Our first main result removes the genericity assumption from the
Baseilhac--Kolb theorem.

\begin{theorem}
\label{thm:intro-root}
Let $\F$ be a field and let $0\neq q\in\F$ be not a root of unity.  For every
total order on the positive roots of $\widehat{\mathfrak{sl}}_2$, the
corresponding ordered monomials in the Baseilhac--Kolb root vectors form a
basis of $\Oq$.  Consequently, Terwilliger's Conjecture~16.1 holds.
\end{theorem}

The proof uses Damiani's explicit straightening relations.  For each fixed
weight, the correction terms strictly decrease a real-root multiplicity
signature, while the leading interchange term has coefficient
$1$, $q^2$, or $q^{-2}$.  A two-level induction on this signature and on the
number of inversions gives the Damiani PBW basis in an arbitrary total order.
The resulting basis is lifted to $\Oq$ through the initial forms of the
Baseilhac--Kolb root vectors.

Our second result concerns the alternating generators.  Put
\[
\begin{array}{lll}
W^-<H<W^+,&W^+<H<W^-,&W^+<W^-<H,\\[1mm]
W^-<W^+<H,&H<W^+<W^-,&H<W^-<W^+.
\end{array}
\tag{1.1}
\label{eq:intro-orders}
\]

\begin{theorem}
\label{thm:intro-alternating}
For $H=G$ and for $H=\widetilde G$, the three corresponding alternating
families give PBW bases of $\Oq$ in each of the six block orders in
\eqref{eq:intro-orders}.  The same twelve bases persist in every $\F$-valued
central fibre of $\Aq$.  In particular, Terwilliger's Conjecture~16.2 holds.
\end{theorem}

The key calculation identifies the initial forms of all alternating
generators with non-zero scalar multiples of the alternating words in the
$q$-shuffle realisation of $\Uplus$.  Terwilliger had previously proved the
corresponding twelve PBW bases in $\Uplus$
\cite{Terwilliger2019Alternating}; filtered lifting then gives
Theorem~\ref{thm:intro-alternating}.  The triangular convolution formula
for the standard tensor factorisation of $\Aq$ shows that arbitrary central
specialisation does not change these initial forms.

For the quotients of Baseilhac and Belliard we obtain the following
consequence.

\begin{corollary}
\label{cor:intro-BB}
Assume that $\operatorname{char}\F\neq2$.  For every non-zero parameter
$\rho$ and every sequence of central values $\{\delta_n\}_{n\geq1}$, the
Baseilhac--Belliard quotient has the conjectured PBW basis in the block order
$W^-<G<W^+$.  Thus their Conjecture~1 holds.
\end{corollary}

Our third result determines two centralisers, in forms stronger than the
original conjectures.

\begin{theorem}
\label{thm:intro-centralisers}
One has
\[
 \Cent_{\Oq}\bigl(\F[W_0,W_{-1},W_{-2},\ldots]\bigr)
 =\F[W_0,W_{-1},W_{-2},\ldots]
\]
and
\[
 \Cent_{\Oq}(\widetilde G_1)
 =\F[\widetilde G_1,\widetilde G_2,\widetilde G_3,\ldots].
\]
Consequently, Terwilliger's Conjectures~16.7 and~16.8 hold.  By symmetry,
the four polynomial subalgebras generated by the four alternating families
are all maximal commutative.
\end{theorem}

The two centraliser arguments are local.  For the negative family we choose
an alternating PBW order in the $q$-shuffle algebra and commute with a
sufficiently high $W_{-N}$.  A positive or imaginary factor creates a unique
maximal profile whose coefficient is a non-zero $q$-geometric sum.  For the
imaginary family we use the Baseilhac--Kolb PBW basis and the action of the
first imaginary root vector on real root vectors.  Again a unique maximal
profile survives with a non-zero geometric coefficient.

The paper is organised as follows.  In Section~2, we define the
$q$-Onsager algebra and its alternating central extension, recall the
word-length filtration, and establish the two filtered lifting lemmas used
throughout.  In Section~3, we introduce the $q$-shuffle algebra, the
alternating words, the Catalan elements, and the Damiani root vectors.  In
Section~4, we prove the arbitrary-order Damiani PBW theorem from explicit
straightening relations.  In Section~5, we compute the initial forms of the
Baseilhac--Kolb root vectors and settle Conjecture~16.1.  In Section~6, we
compute the initial forms of all alternating generators, including the
complete scalar factors and low-degree checks.  In Section~7, we lift the
twelve alternating PBW bases to $\Oq$.  In Section~8, we treat arbitrary
central specialisations and distinguish the three systems of central
coordinates used in the literature.  In Section~9, we prove the conjecture
of Baseilhac and Belliard, first in the normalised case and then for an
arbitrary nonzero parameter.  In Section~10, we establish a large-index
crossing theorem and the negative-tail centraliser.  In Section~11, we
determine the remaining centralisers and deduce maximal commutativity.

\section{The \texorpdfstring{$q$}{q}-Onsager algebra and filtered lifting}
\label{sec:prelim}

Throughout the paper we use the convention
\(\N=\{0,1,2,\ldots\}\).  The symbol \(\F\) denotes a field, and
\(0\neq q\in\F\) is assumed not to be a root of unity.  All algebras are
associative and unital, all algebra homomorphisms preserve the identity,
and all tensor products are taken over \(\F\).  For elements \(X,Y\) of an
\(\F\)-algebra we write
\[
 [n]_q=\frac{q^n-q^{-n}}{q-q^{-1}},\qquad
 \comm{X}{Y}=XY-YX,\qquad
 \qcomm{X}{Y}=qXY-q^{-1}YX .
\]
The assumption on \(q\) will be used repeatedly through the following
elementary observation.

\begin{lemma}\label{lem:nonvanishing-scalars}
For every positive integer \(m\), the scalars
\[
 q-q^{-1},\qquad q+q^{-1},\qquad [m]_q,\qquad
 q^m+q^{-m},\qquad 1-q^{2m},\qquad 1-q^{-2m}
\]
are non-zero.  In particular, the denominators occurring in this section,
and the later geometric factors of the forms \(1-q^{\pm 2m}\) and
\(q^m+q^{-m}\), are non-zero.
\end{lemma}

\begin{proof}
If \(q-q^{-1}=0\), then \(q^2=1\).  If \(q+q^{-1}=0\), then
\(q^2=-1\), and hence \(q^4=1\).  Since \(q-q^{-1}\neq0\), the equality
\([m]_q=0\) implies \(q^m=q^{-m}\), and hence \(q^{2m}=1\).  If
\(q^m+q^{-m}=0\), then \(q^{2m}=-1\), and hence \(q^{4m}=1\).  Finally,
\(1-q^{2m}=0\) or \(1-q^{-2m}=0\) again gives \(q^{2m}=1\).  Each case
contradicts the assumption that \(q\) is not a root of unity.
\end{proof}

\subsection{The \texorpdfstring{\(q\)}{q}-Onsager algebras and the word-length filtration}

\begin{definition}
\label{def:Oq}
The \(q\)-Onsager algebra \(\Oq\) is the \(\F\)-algebra generated by
\(W_0,W_1\), subject to the \(q\)-Dolan--Grady relations
\begin{align}
W_0^3W_1-[3]_qW_0^2W_1W_0+[3]_qW_0W_1W_0^2-W_1W_0^3
   &=(q^2-q^{-2})^2(W_1W_0-W_0W_1),
\label{eq:qDG0}\\
W_1^3W_0-[3]_qW_1^2W_0W_1+[3]_qW_1W_0W_1^2-W_0W_1^3
   &=(q^2-q^{-2})^2(W_0W_1-W_1W_0).
\label{eq:qDG1}
\end{align}
\end{definition}

%Our normalisation agrees with \cite{Terwilliger2022OqACE}.  
Equivalently,
the right-hand sides of \eqref{eq:qDG0} and \eqref{eq:qDG1} are
\(\rho_0\comm{W_0}{W_1}\) and \(\rho_0\comm{W_1}{W_0}\), respectively, where
\begin{equation}
\label{eq:rho0}
        \rho_0=-(q^2-q^{-2})^2 .
\end{equation}

We shall also use Terwilliger's alternating central extension of
\(\Oq\).  The current algebra introduced by Baseilhac and Koizumi and
presented by Baseilhac and Shigechi is isomorphic to this algebra
\cite{BaseilhacKoizumi2005,BaseilhacShigechi2010,Terwilliger2021ACE}.  We
use the following presentation as its definition.

\begin{definition}\label{def:Aq}
The alternating central extension $\Aq$ of the $q$-Onsager algebra is the
$\F$-algebra generated by the four families
\[
        \cW_{-k},\qquad \cW_{k+1},\qquad
        \cG_{k+1},\qquad \Gt_{k+1}
        \qquad (k\in\N).
\]
These generators are called the \textbf{alternating generators} of $\Aq$.
The defining relations are, for all $k,\ell\in\N$:
\begin{align}
&[\cW_0,\cW_{k+1}]
        =[\cW_{-k},\cW_1]
          =\frac{\Gt_{k+1}-\cG_{k+1}}{q+q^{-1}},
          \label{eq:Aq1}\\
&[\cW_0,\cG_{k+1}]_q
        =[\Gt_{k+1},\cW_0]_q
          =\rho \cW_{-k-1}-\rho \cW_{k+1},
          \label{eq:Aq2}\\
&[\cG_{k+1},\cW_1]_q
        =[\cW_1,\Gt_{k+1}]_q
          =\rho \cW_{k+2}-\rho \cW_{-k},
          \label{eq:Aq3}\\
&[\cW_{-k},\cW_{-\ell}]=0,
        \qquad
        [\cW_{k+1},\cW_{\ell+1}]=0,
        \label{eq:Aq4}\\
&[\cW_{-k},\cW_{\ell+1}]
        +[\cW_{k+1},\cW_{-\ell}]=0,
        \label{eq:Aq5}\\
&[\cW_{-k},\cG_{\ell+1}]
        +[\cG_{k+1},\cW_{-\ell}]=0,
        \label{eq:Aq6}\\
&[\cW_{-k},\Gt_{\ell+1}]
        +[\Gt_{k+1},\cW_{-\ell}]=0,
        \label{eq:Aq7}\\
&[\cW_{k+1},\cG_{\ell+1}]
        +[\cG_{k+1},\cW_{\ell+1}]=0,
        \label{eq:Aq8}\\
&[\cW_{k+1},\Gt_{\ell+1}]
        +[\Gt_{k+1},\cW_{\ell+1}]=0,
        \label{eq:Aq9}\\
&[\cG_{k+1},\cG_{\ell+1}]=0,
        \qquad
        [\Gt_{k+1},\Gt_{\ell+1}]=0,
        \label{eq:Aq10}\\
&[\Gt_{k+1},\cG_{\ell+1}]
        +[\cG_{k+1},\Gt_{\ell+1}]=0,
        \label{eq:Aq11}
\end{align}
where $\rho=-(q^2-q^{-2})^2$.
\end{definition}

We use the scalar convention
\begin{equation}
\label{eq:Aq-G0}
        \cG_0=\ctG_0=-(q-q^{-1})[2]_q^2 .
\end{equation}
These two elements are scalars, not additional generators.

Terwilliger~\cite{Terwilliger2021ACE,Terwilliger2022OqACE} proved that the centre \(Z(\Aq)\) is a polynomial algebra in
countably many variables and that the standard multiplication map gives a
tensor product factorisation
\begin{equation}
\label{eq:standard-factorisation-intro}
        \Aq\simeq \Oq\otimes\F[z_1,z_2,\ldots].
\end{equation}
Under this factorisation, evaluation of the polynomial variables at
\(z_1=z_2=\cdots=0\) gives the canonical central reduction
\[
        \gamma:\Aq\longrightarrow\Oq .
\]
The alternating generators of \(\Oq\) are, by definition, the images
\[
 W_{-n}=\gamma(\cW_{-n}),\qquad
 W_{n+1}=\gamma(\cW_{n+1}),\qquad
 G_{n+1}=\gamma(\cG_{n+1}),\qquad
 \widetilde G_{n+1}=\gamma(\ctG_{n+1})
 \qquad(n\in\N).
\]
Thus calligraphic alternating generators belong to \(\Aq\), while the
corresponding plain symbols denote their canonical images in \(\Oq\).

We next recall the filtration that connects \(\Oq\) with the positive part
of \(U_q(\widehat{\mathfrak{sl}}_2)\).  For \(d\in\N\), let \(F_d\Oq\) be
the \(\F\)-subspace spanned by the images of all words in \(W_0,W_1\) of
length at most \(d\), and put \(F_{-1}\Oq=0\).  Then
\(F_r\Oq\,F_s\Oq\subseteq F_{r+s}\Oq\), and
\(\Oq=\bigcup_{d\geq0}F_d\Oq\).  We call this the \textbf{word-length filtration}.
The associated graded algebra is
\[
        \gr\Oq=\bigoplus_{d\geq0}F_d\Oq/F_{d-1}\Oq .
\]
If \(u\in F_d\Oq\), we write \(\overline u\) for its \textbf{image} in
\(F_d\Oq/F_{d-1}\Oq\).  If \(u\in F_d\Oq\setminus F_{d-1}\Oq\), then
\(\overline u\) is called the \textbf{initial form} of \(u\), and is denoted by
\(\init(u)\).

Let \(\Uplus\) be the \(\F\)-algebra generated by \(A,B\), subject to the
cubic \(q\)-Serre relations
\begin{align*}
A^3B-[3]_qA^2BA+[3]_qABA^2-BA^3&=0,\\
B^3A-[3]_qB^2AB+[3]_qBAB^2-AB^3&=0.
\end{align*}
The degree-four parts of the \(q\)-Dolan--Grady relations are precisely
these two relations.  The following theorem makes this degeneration exact.

\begin{theorem}[{{\cite[Theorem~4.4]{Terwilliger2017PositivePart}}}]
\label{thm:gr-isomorphism}
There is an algebra isomorphism
\[
        \psi:\Uplus\longrightarrow\gr\Oq
\]
such that \(\psi(A)=\overline{W_0}\) and
\(\psi(B)=\overline{W_1}\).
\end{theorem}

From now on, we identify \(\gr\Oq\) with \(\Uplus\) through \(\psi\).

\subsection{Filtered lifting, centralisers, and notation}

We use the following convention for PBW bases.

\begin{definition}
Let \(R\) be an \(\F\)-algebra, let \(\Omega\subseteq R\), and let \(<\) be
a linear order on \(\Omega\).  We say that \(\Omega\), with the order \(<\),
gives a PBW basis of \(R\) if the ordered monomials
\[
        a_1a_2\cdots a_m,
        \qquad
        m\in\N,\qquad
        a_1\leq a_2\leq\cdots\leq a_m,
\]
form an \(\F\)-basis of \(R\).  For \(m=0\), the monomial is interpreted as
the identity element.
\end{definition}

The next lemma is the filtered lifting principle used throughout the paper.
It is stated in the positive-degree form needed below.

\begin{lemma}
\label{lem:filtered-lifting}
Let \(R\) be an \(\F\)-algebra with an exhaustive filtration
\[
        0=F_{-1}R\subseteq F_0R\subseteq F_1R\subseteq\cdots .
\]
Let \(\{u_\lambda\}_{\lambda\in\Lambda}\) be elements of positive filtration
degree, say
\[
        u_\lambda\in F_{d_\lambda}R\setminus F_{d_\lambda-1}R,
        \qquad d_\lambda>0.
\]
Assume that
\[
        \init(u_\lambda)=c_\lambda v_\lambda,
        \qquad c_\lambda\in\F^\times,
\]
where \(v_\lambda\in\gr R\) is homogeneous of degree \(d_\lambda\).  Fix a
linear order on \(\Lambda\).  If the ordered monomials in the
\(v_\lambda\) form a basis of \(\gr R\), then the corresponding ordered
monomials in the \(u_\lambda\) form a basis of \(R\).
\end{lemma}

\begin{proof}
For an ordered monomial \(M\), write \(M(u)\) and \(M(v)\) for the
corresponding products in the \(u_\lambda\) and in the \(v_\lambda\),
respectively.  If \(M\) contains \(u_\lambda\) with multiplicity
\(e_\lambda(M)\), then the initial form of \(M(u)\) is
\[
        \left(\prod_\lambda c_\lambda^{e_\lambda(M)}\right)M(v).
\]
In particular, \(M(u)\) has the expected filtration degree, because
\(M(v)\) is a non-zero basis element of \(\gr R\).

We first prove linear independence.  Suppose that there is a non-trivial
finite relation
\[
        \sum_M a_M M(u)=0
\]
among distinct ordered monomials.  Let \(D\) be the largest filtration
degree of a monomial occurring with non-zero coefficient.  Passing to
\(F_DR/F_{D-1}R\) gives
\[
 \sum_{\deg M=D}
 a_M\left(\prod_\lambda c_\lambda^{e_\lambda(M)}\right)M(v)=0 .
\]
The monomials \(M(v)\) are distinct basis elements of \(\gr R\), and all
displayed scalar factors are non-zero.  Hence \(a_M=0\) for every monomial
of degree \(D\).  Repeating this argument in descending filtration degree
gives \(a_M=0\) for all \(M\), a contradiction.

We next prove spanning by induction on the filtration degree.  Let
\(r\in F_DR\).  The image of \(r\) in \(F_DR/F_{D-1}R\) is a finite linear
combination of ordered monomials \(M(v)\) of degree \(D\).  Replacing each
\(M(v)\) by
\[
        \left(\prod_\lambda c_\lambda^{e_\lambda(M)}\right)^{-1}M(u)
\]
gives an element \(r'\in F_DR\) with the same image as \(r\) in
\(F_DR/F_{D-1}R\).  Thus \(r-r'\in F_{D-1}R\), and the induction hypothesis
completes the proof.
\end{proof}

\begin{remark}
Lemma~\ref{lem:filtered-lifting} is an initial-form version of
\cite[Proposition~4.6]{Terwilliger2017PositivePart}.  The formulation above
does not require the lifting elements to be represented in advance by
homogeneous words in a free algebra.
\end{remark}

We shall also use the corresponding lifting statement for centralisers.  If
\(S\) is a subset of an algebra \(R\), we write \(\Cent_R(S)\) for its
centraliser.

\begin{lemma}\label{lem:filtered-centraliser}
Let \(R=\bigcup_{d\geq0}F_dR\) be an exhaustively filtered algebra, with
\(F_{-1}R=0\), and let \(H\subseteq R\) be a commutative subalgebra filtered
by \(H\cap F_dR\).  View \(\gr H\) as the corresponding graded subalgebra of
\(\gr R\).  If
\[
        \Cent_{\gr R}(\gr H)=\gr H,
\]
then
\[
        \Cent_R(H)=H .
\]
\end{lemma}

\begin{proof}
Since \(H\) is commutative, \(H\subseteq\Cent_R(H)\).  Conversely, let
\(0\neq x\in\Cent_R(H)\), and let \(D\) be the filtration degree of \(x\).
For every homogeneous element of \(\gr H\), choose a representative
\(h\in H\).  The equality \([x,h]=0\) implies that the highest-degree
component of \([x,h]\) is zero; hence \(\init(x)\) centralises every
homogeneous element of \(\gr H\).  Therefore
\[
        \init(x)\in \Cent_{\gr R}(\gr H)=\gr H .
\]
By the definition of \(\gr H\), there exists \(y\in H\cap F_DR\) with
\(\init(y)=\init(x)\).  Then \(x-y\in F_{D-1}R\), and \(x-y\) still
centralises \(H\).  Induction on \(D\) gives \(x-y\in H\), and hence
\(x\in H\).
\end{proof}

We shall use three standard symmetries of \(\Oq\).  There is an involutive
automorphism \(\sigma\), an involutive antiautomorphism \(\dagger\), and an
involutive antiautomorphism \(\tau\), satisfying, for \(n\in\N\),
\begin{align*}
\sigma:&\quad
 W_{-n}\leftrightarrow W_{n+1},
 \qquad
 G_{n+1}\leftrightarrow\widetilde G_{n+1},\\
\dagger:&\quad
 W_{-n}\mapsto W_{-n},
 \qquad
 W_{n+1}\mapsto W_{n+1},
 \qquad
 G_{n+1}\leftrightarrow\widetilde G_{n+1},\\
\tau:&\quad
 W_{-n}\leftrightarrow W_{n+1},
 \qquad
 G_{n+1}\mapsto G_{n+1},
 \qquad
 \widetilde G_{n+1}\mapsto\widetilde G_{n+1}.
\end{align*}
See \cite[Definitions~3.4--3.6 and Theorem~11.6]{Terwilliger2022OqACE}.
These maps send \(W_0,W_1\) to elements of filtration degree one, and
therefore preserve the word-length filtration.

\paragraph{Typographical conventions.}
Calligraphic symbols such as
\(\cW_{-n},\cW_{n+1},\cG_n,\ctG_n\) always denote elements of \(\Aq\).
The corresponding plain symbols
\(W_{-n},W_{n+1},G_n,\widetilde G_n\) denote their canonical images in
\(\Oq\).  From Section~\ref{sec:qshuffle} onwards, the superscript
\(\sh\) is reserved for elements of the \(q\)-shuffle algebra.  Thus the
notation distinguishes, for example, \(\cW_{-n}\in\Aq\),
\(W_{-n}\in\Oq\), and \(W_{-n}^{\sh}\) in the \(q\)-shuffle algebra.




\section{The \texorpdfstring{$q$}{q}-shuffle algebra, alternating words, and root vectors}
\label{sec:qshuffle}

In this section we recall the \(q\)-shuffle realisation of \(\Uplus\), the
alternating words, the Catalan elements, and Damiani's root vectors.  These
objects will provide the associated-graded structure for the PBW lifting
arguments in later sections.

\subsection{The \texorpdfstring{\(q\)}{q}-shuffle realisation and the alternating words}

Let \(\Vsh\) be the \(\F\)-vector space with basis the set of all words in
the two letters \(x,y\), including the empty word \(1\).  Juxtaposition of
letters and words always denotes concatenation.  The \(q\)-shuffle product
will be denoted by \(\star\).

Define a symmetric bilinear form on the letters by
\[
 \langle x,x\rangle=\langle y,y\rangle=2,\qquad
 \langle x,y\rangle=\langle y,x\rangle=-2 .
\]
The \textbf{\(q\)-shuffle product} is the unique bilinear product on \(\Vsh\)
determined by \(1\star v=v\star1=v\) and by the recursion
\begin{equation}
\label{eq:qshuffle-recursion}
\begin{split}
 u\star v
={}&
 u_1\bigl((u_2\cdots u_r)\star v\bigr)+
 v_1\bigl(u\star(v_2\cdots v_s)\bigr)
 q^{\langle u_1,v_1\rangle+\cdots+\langle u_r,v_1\rangle},
\end{split}
\end{equation}
for non-empty words \(u=u_1\cdots u_r\) and \(v=v_1\cdots v_s\).  This
product is associative; see \cite[Section~4, equations~(8)--(12)]
{Terwilliger2019Alternating}.  For instance,
\[
        x\star y=xy+q^{-2}yx,\qquad
        y\star x=yx+q^{-2}xy .
\]

Rosso's \(q\)-shuffle embedding \cite{Rosso1998} gives an injective algebra homomorphism
\[
        \iota:\Uplus\longrightarrow(\Vsh,\star),
        \qquad
        A\longmapsto x,\qquad B\longmapsto y .
\]
By \cite[Definition~4.2 and the paragraph following equations~(14)--(15)]{Terwilliger2019Alternating}, the image is the
subalgebra of \((\Vsh,\star)\) generated by \(x,y\).  
We identify
\(\Uplus\) with this image.  Combined with
Theorem~\ref{thm:gr-isomorphism}, this identification lets us regard
initial forms in \(\gr\Oq\) as elements of the \(q\)-shuffle algebra.

For \(n\in\N\), define the alternating words
\begin{equation}
\label{eq:alternating-words}
 W_{-n}^{\sh}=x(yx)^n,\qquad
 W_{n+1}^{\sh}=y(xy)^n,\qquad
 G_n^{\sh}=(yx)^n,\qquad
 \widetilde G_n^{\sh}=(xy)^n .
\end{equation}
Here the powers are concatenation powers of words.  Thus
\[
 W_0^{\sh}=x,\qquad W_1^{\sh}=y,\qquad
 G_0^{\sh}=\widetilde G_0^{\sh}=1,\qquad
 G_1^{\sh}=yx,\qquad \widetilde G_1^{\sh}=xy .
\]
These are precisely the alternating words with a superscript
added to distinguish them from alternating generators in \(\Oq\); see~\cite[Definitions~5.1--5.2]{Terwilliger2019Alternating}.

The word-reversal map fixes each of \(x\) and \(y\).  As a map on
\((\Vsh,\star)\), it is an antiautomorphism; see
\cite[Lemma~4.1]{Terwilliger2019Alternating}.  It fixes
\(W_{-n}^{\sh}\) and \(W_{n+1}^{\sh}\), and interchanges
\(G_n^{\sh}\) with \(\widetilde G_n^{\sh}\).

Whenever a commutator or \(q\)-commutator is written inside \(\Vsh\), the
products are computed using \(\star\).  We shall use the following two
relations.  For \(n\geq1\),
\begin{align}
\qcomm{\widetilde G_n^{\sh}}{x}
 &=(q-q^{-1})W_{-n}^{\sh},
\label{eq:qshuffle-Wminus}\\
\qcomm{y}{\widetilde G_n^{\sh}}
 &=(q-q^{-1})W_{n+1}^{\sh}.
\label{eq:qshuffle-Wplus}
\end{align}
They are the cases \(k=n-1\) of
\cite[Proposition~5.7, equations~(38)--(39)]
{Terwilliger2019Alternating}.

The following theorem is the graded PBW structure used throughout the paper.

\begin{theorem}[{{\cite[Theorems~10.1 and~10.2]{Terwilliger2019Alternating}}}]
\label{thm:qshuffle-twelve}
The three families
\[
 \{W_{-i}^{\sh}\}_{i\in\N},\qquad
 \{\widetilde G_{j+1}^{\sh}\}_{j\in\N},\qquad
 \{W_{k+1}^{\sh}\}_{k\in\N}
\]
give a PBW basis of \(\Uplus\) in every linear order satisfying one of the
six block conditions
\begin{equation}
\label{eq:six-shuffle-orders}
\begin{array}{lll}
W^-<\widetilde G<W^+,&
W^+<\widetilde G<W^-,&
W^+<W^-<\widetilde G,\\[1mm]
W^-<W^+<\widetilde G,&
\widetilde G<W^+<W^-,&
\widetilde G<W^-<W^+.
\end{array}
\end{equation}
No condition is imposed on the order inside a single block.  The same
assertion holds with the family
\(\{\widetilde G_{j+1}^{\sh}\}_{j\in\N}\) replaced by
\(\{G_{j+1}^{\sh}\}_{j\in\N}\).
\end{theorem}


\subsection{Catalan elements and Damiani root vectors}

We next recall the Catalan elements that describe Damiani's root vectors in
the \(q\)-shuffle realisation.  Let
\[
        \epsilon(x)=1,\qquad \epsilon(y)=-1 .
\]
A word \(v=v_1\cdots v_r\) in the letters \(x,y\) is called Catalan if
\[
 \sum_{h=1}^{i}\epsilon(v_h)\geq0\quad(1\leq i<r),
 \qquad
 \sum_{h=1}^{r}\epsilon(v_h)=0 .
\]
Such a word has even length.  For \(n\in\N\), define
\begin{equation}
\label{eq:Catalan-definition}
 C_n=
 \sum_{v=v_1\cdots v_{2n}}
 v\prod_{j=0}^{2n}
 \left[1+\sum_{h=1}^{j}\epsilon(v_h)\right]_q ,
\end{equation}
where the sum is over all Catalan words of length \(2n\), and the inner sum
is interpreted as \(0\) when \(j=0\).  Thus
\[
        C_0=1,\qquad C_1=[2]_qxy .
\]
This is the Catalan element normalisation used in
\cite[Section~4, equations following Definition~4.2]
{Terwilliger2019Alternating}; equivalently, it is the normalisation from
\cite[Theorem~1.7]{Terwilliger2019Catalan}.

Let
\[
 \{E_{n\delta+\alpha_0}\}_{n\geq0},\qquad
 \{E_{n\delta+\alpha_1}\}_{n\geq0},\qquad
 \{E_{n\delta}\}_{n\geq1}
\]
be Damiani's root vectors in \(\Uplus\), in the normalisation
\cite[Section~2, equations~(4)--(6)]{Terwilliger2019Alternating}.  They are
defined recursively by
\begin{align}
&E_{\alpha_0}=A,
\qquad E_{\alpha_1}=B,
\qquad E_\delta=q^{-2}BA-AB,
\label{eq:Damiani-initial}\\
&E_{n\delta+\alpha_0}
 =\frac{\comm{E_\delta}{E_{(n-1)\delta+\alpha_0}}}{[2]_q},
\qquad E_{n\delta+\alpha_1}
 =\frac{\comm{E_{(n-1)\delta+\alpha_1}}{E_\delta}}{[2]_q},
\label{eq:Damiani-real}\\
&E_{n\delta}
 =q^{-2}E_{(n-1)\delta+\alpha_1}A
   -AE_{(n-1)\delta+\alpha_1}
\qquad(n\geq1).
\label{eq:Damiani-imaginary}
\end{align}

Under the \(q\)-shuffle embedding \(\iota\), these root vectors have the
closed forms
\begin{align}
E_{n\delta+\alpha_0}
 &\longmapsto q^{-2n}(q-q^{-1})^{2n}xC_n
 \qquad(n\geq0),
\label{eq:Damiani-shuffle-0}\\
E_{n\delta+\alpha_1}
 &\longmapsto q^{-2n}(q-q^{-1})^{2n}C_ny
 \qquad(n\geq0),
\label{eq:Damiani-shuffle-1}\\
E_{n\delta}
 &\longmapsto -q^{-2n}(q-q^{-1})^{2n-1}C_n
 \qquad(n\geq1).
\label{eq:Damiani-shuffle-delta}
\end{align}
Here \(xC_n\) and \(C_ny\) denote concatenation of words, not
\(q\)-shuffle multiplication.  The formulae are
\cite[Theorem~1.7]{Terwilliger2019Catalan}; in the present notation they
are also recorded in
\cite[Section~4, after Definition~4.2]{Terwilliger2019Alternating}.

Finally, we record the Catalan--alternating relation that will be used to
compute the highest filtered component of the alternating generators of
\(\Oq\).

\begin{theorem}[{{\cite[Theorem~11.14]{Terwilliger2019Alternating}}}]
\label{thm:Catalan-alternating}
For every \(n\in\N\),
\begin{equation}
\label{eq:Catalan-alternating}
        \sum_{i=0}^{n}(-1)^i[2n-i]_q\,
        C_i\star\widetilde G_{n-i}^{\sh}=0 .
\end{equation}
\end{theorem}



\section{Root-vector PBW bases away from roots of unity}
\label{sec:specialisation}

In this section we prove that Damiani's root vectors give a PBW basis in an
arbitrary total order whenever $q$ is not a root of unity.  The point is
that the explicit straightening relations have leading coefficients among
$1,q^2,q^{-2}$, while every correction term decreases a fixed finite
triangular parameter.  Thus no generic-parameter or localisation argument
is needed.

For $i\in\N$ and $m\geq1$, abbreviate
\[
 A_i=E_{i\delta+\alpha_0},\qquad
 B_i=E_{i\delta+\alpha_1},\qquad
 I_m=E_{m\delta}.
\]
We use the following standard monomial convention:
\begin{equation}
\label{eq:Damiani-standard-order}
 A_i<A_{i+1},\qquad I_m<I_{m+1},\qquad B_{j+1}<B_j,
 \qquad A_i<I_m<B_j
\end{equation}
for all $i,j\in\N$ and $m\geq1$.  Thus the imaginary root vectors occur in
increasing index order.  The standard ordered monomials in these root
vectors form a basis of $\Uplus$; see
\cite[Proposition~2.4]{Terwilliger2019Alternating} and
\cite[Section~5, Theorem~2]{Damiani1993}.  Baseilhac--Kolb use the opposite
internal order on the imaginary string.  Since the elements $I_m$ commute
pairwise, the two conventions give the same set of standard monomials.

For a word $w$ in the root vectors, let $a_i(w)$ and $b_i(w)$ denote the
multiplicities of $A_i$ and $B_i$, respectively, and define its
\emph{real signature} by
\begin{equation}
\label{eq:real-signature}
 \operatorname{rsig}(w)
 =\bigl(a_0(w),a_1(w),a_2(w),\ldots;
        b_0(w),b_1(w),b_2(w),\ldots\bigr).
\end{equation}
Only finitely many entries are non-zero.  We order these signatures
lexicographically, reading the $A$-coordinates first and then the
$B$-coordinates, and declaring the signature with the larger entry at the
first place of disagreement to be larger.  In a fixed weight there are only
finitely many signatures.  Moreover, strict comparison is preserved by
adding the same multiplicity vector:
\begin{equation}
\label{eq:signature-additivity}
 \boldsymbol a<\boldsymbol b
 \quad\Longrightarrow\quad
 \boldsymbol a+\boldsymbol c<\boldsymbol b+\boldsymbol c.
\end{equation}

\begin{lemma}[Triangular Damiani crossings]
\label{lem:Damiani-triangular-crossings}
Let $X,Y$ be two distinct Damiani root vectors.  There exist
$\varepsilon_{X,Y}\in\{1,q^2,q^{-2}\}$ and finitely many scalars
$c_s\in\F$ and words $w_s$ in Damiani root vectors such that
\begin{equation}
\label{eq:Damiani-triangular-crossing}
 XY=\varepsilon_{X,Y}YX+\sum_s c_sw_s,
\end{equation}
where every $w_s$ has the same weight as $XY$ and satisfies
\[
 \operatorname{rsig}(w_s)<\operatorname{rsig}(XY).
\]
The same assertion holds when either orientation of a crossing relation is
solved for the other product.
\end{lemma}

\begin{proof}
We use the explicit relations in
\cite[Lemmas~3.1--3.3]{Terwilliger2019Catalan}.  In the notation above they
are as follows.  For $i,j\in\N$ and $m\geq1$,
\begin{align}
 B_iA_j&=q^2A_jB_i+q^2I_{i+j+1},
 \label{eq:Damiani-cross-BA}\\
 I_mA_j&=A_jI_m+q^{2-2m}[2]_qA_{m+j}
 -q^2(q^2-q^{-2})\sum_{\ell=1}^{m-1}
 q^{-2\ell}A_{j+\ell}I_{m-\ell},
 \label{eq:Damiani-cross-IA}\\
 B_jI_m&=I_mB_j+q^{2-2m}[2]_qB_{m+j}
 -q^2(q^2-q^{-2})\sum_{\ell=1}^{m-1}
 q^{-2\ell}I_{m-\ell}B_{j+\ell}.
 \label{eq:Damiani-cross-BI}
\end{align}
For $i>j\geq0$, if $i-j=2r+1$ is odd, then
\begin{align}
 A_iA_j&=q^{-2}A_jA_i
 -(q^2-q^{-2})\sum_{\ell=1}^{r}q^{-2\ell}
 A_{j+\ell}A_{i-\ell},
 \label{eq:Damiani-cross-AA-odd}\\
 B_jB_i&=q^{-2}B_iB_j
 -(q^2-q^{-2})\sum_{\ell=1}^{r}q^{-2\ell}
 B_{i-\ell}B_{j+\ell}.
 \label{eq:Damiani-cross-BB-odd}
\end{align}
If $i-j=2r$ is even, then
\begin{align}
 A_iA_j={}&q^{-2}A_jA_i
 -q^{j-i+1}(q-q^{-1})A_{j+r}^{2}
 \notag\\
 &-(q^2-q^{-2})\sum_{\ell=1}^{r-1}q^{-2\ell}
 A_{j+\ell}A_{i-\ell},
 \label{eq:Damiani-cross-AA-even}\\
 B_jB_i={}&q^{-2}B_iB_j
 -q^{j-i+1}(q-q^{-1})B_{j+r}^{2}
 \notag\\
 &-(q^2-q^{-2})\sum_{\ell=1}^{r-1}q^{-2\ell}
 B_{i-\ell}B_{j+\ell}.
 \label{eq:Damiani-cross-BB-even}
\end{align}
Empty sums are zero, and $I_mI_n=I_nI_m$ for all $m,n\geq1$.

We check triangularity.  In \eqref{eq:Damiani-cross-BA}, the correction
removes one $A$-factor and one $B$-factor.  In
\eqref{eq:Damiani-cross-IA} and \eqref{eq:Damiani-cross-BI}, each
correction replaces a real root of index $j$ by one of strictly larger
index.  Thus the first real multiplicity that changes decreases.  In
\eqref{eq:Damiani-cross-AA-odd}--\eqref{eq:Damiani-cross-BB-even}, each
correction replaces the pair of numerical indices $j<i$ by indices
strictly larger than $j$; in the even case the midpoint occurs twice.
Again the multiplicity at the first changed coordinate, namely index $j$,
decreases.  Therefore every correction word has strictly smaller real
signature.

The leading coefficients are $1$, $q^2$, or $q^{-2}$, and are non-zero.
Consequently each relation may be solved in the reverse direction.  This
only rescales the same correction words, so their strict signature decrease
is preserved.
\end{proof}

\begin{theorem}
\label{thm:any-Damiani}
Let $\prec$ be any total order on the set $\Phi_+$ of positive roots of
$\widehat{\mathfrak{sl}}_2$.  The $\prec$-ordered monomials in Damiani's
root vectors form an $\F$-basis of $U_q^+$.
\end{theorem}

\begin{proof}
Fix a weight $\nu$.  Only finitely many positive roots and finitely many
root partitions occur in weight $\nu$.  For a word $w=X_1\cdots X_d$ of
this weight, let $\operatorname{inv}_\prec(w)$ be the number of pairs
$r<s$ for which $X_s\prec X_r$.

We prove, by induction on the real signature and then on
$\operatorname{inv}_\prec$, that every word of weight $\nu$ is a linear
combination of $\prec$-ordered monomials.  If a word is not ordered, choose
an adjacent $\prec$-inversion $XY$.  Apply
Lemma~\ref{lem:Damiani-triangular-crossings}, using the displayed relation
or its reverse according to the Damiani order.  The leading term interchanges
$X$ and $Y$; it has the same real signature and one fewer
$\prec$-inversion.  Each correction term has smaller real signature, even
after the unchanged factors on the left and right are restored, by
\eqref{eq:signature-additivity}.  The two induction hypotheses therefore
handle every term.  Since the set of signatures in weight $\nu$ is finite,
the induction terminates.  Thus the $\prec$-ordered monomials span the
$\nu$-weight space.

For every root partition of $\nu$ there is exactly one $\prec$-ordered
monomial, so this spanning set has the same cardinality as the standard
Damiani PBW basis in weight $\nu$.  The latter is a basis by Damiani's PBW
theorem, and hence the weight space has that dimension.  The
$\prec$-ordered spanning set is therefore a basis.  Taking the direct sum
over all weights proves the theorem.
\end{proof}


\section{Initial forms of the Baseilhac--Kolb root vectors}
\label{sec:root-initial}

In this section we compute the initial forms of the Baseilhac--Kolb root
vectors in \(\Oq\).  These computations are the bridge between the
Baseilhac--Kolb root-vector construction and Damiani's root vectors in
\(\Uplus\simeq\gr\Oq\).  They will also be used in
Section~\ref{sec:alternating-initial} to connect the recursion for the
alternating generators with the Catalan relation
Theorem~\ref{thm:Catalan-alternating}.

For a non-zero element \(u\in\Oq\), write
\[
        \deg_F u=\min\{d\mid u\in F_d\Oq\}
\]
for its word-length filtration degree.  We continue to identify
\(\gr\Oq\) with \(\Uplus\) through Theorem~\ref{thm:gr-isomorphism}.

The Baseilhac--Kolb root vectors in \(\Oq\) are denoted
\[
 \{B_{n\delta+\alpha_0}\}_{n\geq0},\qquad
 \{B_{n\delta+\alpha_1}\}_{n\geq0},\qquad
 \{B_{n\delta}\}_{n\geq1}.
\]
In the normalisation used here, they are recursively defined as follows.
First,
\begin{equation}
\label{eq:Bdelta}
        B_\delta=q^{-2}W_1W_0-W_0W_1 .
\end{equation}
The real root vectors are given by
\begin{align}
B_{\alpha_0}&=W_0,
&
B_{\delta+\alpha_0}
 &=W_1+
   \frac{q\,\comm{B_\delta}{W_0}}
        {(q-q^{-1})(q^2-q^{-2})},
\label{eq:BK-real0-first}\\
B_{n\delta+\alpha_0}
 &=B_{(n-2)\delta+\alpha_0}
   +\frac{q\,\comm{B_\delta}{B_{(n-1)\delta+\alpha_0}}}
        {(q-q^{-1})(q^2-q^{-2})}
        \qquad(n\geq2),
\label{eq:BK-real0}\\
B_{\alpha_1}&=W_1,
&
B_{\delta+\alpha_1}
 &=W_0-
   \frac{q\,\comm{B_\delta}{W_1}}
        {(q-q^{-1})(q^2-q^{-2})},
\label{eq:BK-real1-first}\\
B_{n\delta+\alpha_1}
 &=B_{(n-2)\delta+\alpha_1}
   -\frac{q\,\comm{B_\delta}{B_{(n-1)\delta+\alpha_1}}}
        {(q-q^{-1})(q^2-q^{-2})}
        \qquad(n\geq2).
\label{eq:BK-real1}
\end{align}
For \(n\geq1\), the imaginary root vectors are given by
\begin{align}
B_{n\delta}
={}&q^{-2}B_{(n-1)\delta+\alpha_1}W_0
    -W_0B_{(n-1)\delta+\alpha_1}+(q^{-2}-1)
  \sum_{\ell=0}^{n-2}
  B_{\ell\delta+\alpha_1}
  B_{(n-\ell-2)\delta+\alpha_1}.
\label{eq:BK-imaginary}
\end{align}
An empty sum is interpreted as zero.  These formulae are recalled in
\cite[Section~4, equations~(7)--(12)]{Terwilliger2022OqACE}; they come
from the Baseilhac--Kolb construction of root vectors in
\cite[Section~3]{BaseilhacKolb2020}.

We shall repeatedly use the identity
\begin{equation}
\label{eq:denominator-expanded}
 (q-q^{-1})(q^2-q^{-2})=(q-q^{-1})^2[2]_q .
\end{equation}
By Lemma~\ref{lem:nonvanishing-scalars}, all denominators in
\eqref{eq:BK-real0-first}--\eqref{eq:BK-real1} are non-zero.

\begin{proposition}
\label{prop:root-leading-U}
For \(n\in\N\), the real root vectors have initial forms
\begin{align}
\init(B_{n\delta+\alpha_0})
 &=q^n(q-q^{-1})^{-2n}E_{n\delta+\alpha_0},
\label{eq:root-leading-real0}\\
\init(B_{n\delta+\alpha_1})
 &=q^n(q-q^{-1})^{-2n}E_{n\delta+\alpha_1}.
\label{eq:root-leading-real1}
\end{align}
For \(n\geq1\), the imaginary root vectors have initial forms
\begin{equation}
\label{eq:root-leading-imaginary}
        \init(B_{n\delta})
        =q^{\,n-1}(q-q^{-1})^{2-2n}E_{n\delta}.
\end{equation}
In particular,
\[
        \deg_F B_{n\delta+\alpha_0}
        =\deg_F B_{n\delta+\alpha_1}=2n+1,
        \qquad
        \deg_F B_{n\delta}=2n .
\]
\end{proposition}

\begin{proof}
By \eqref{eq:Bdelta} and Theorem~\ref{thm:gr-isomorphism},
\[
        \init(B_\delta)=q^{-2}BA-AB=E_\delta .
\]

We first prove \eqref{eq:root-leading-real0}.  For \(n=0\), this is just
\(\init(B_{\alpha_0})=\init(W_0)=A=E_{\alpha_0}\).

Consider \(n=1\).  By \eqref{eq:denominator-expanded}, the coefficient of
\(\comm{B_\delta}{W_0}\) in \eqref{eq:BK-real0-first} is
\[
        \frac{q}{(q-q^{-1})^2[2]_q}.
\]
The summand \(W_1\) has filtration degree one.  The commutator
\(\comm{B_\delta}{W_0}\) has filtration degree at most three, and its
degree-three component is \(\comm{E_\delta}{A}\).  Hence the degree-three
component of \(B_{\delta+\alpha_0}\) is
\[
        \frac{q}{(q-q^{-1})^2[2]_q}\,
        \comm{E_\delta}{A}
        =
        q(q-q^{-1})^{-2}E_{\delta+\alpha_0},
\]
where the last equality uses \eqref{eq:Damiani-real}.  This is non-zero, so
\eqref{eq:root-leading-real0} holds for \(n=1\).

Now let \(n\geq2\), and assume that
\[
        \init(B_{(n-1)\delta+\alpha_0})
        =
        q^{n-1}(q-q^{-1})^{-2n+2}
        E_{(n-1)\delta+\alpha_0}.
\]
The first term on the right-hand side of \eqref{eq:BK-real0} has filtration
degree \(2n-3\).  The commutator term has filtration degree at most
\(2n+1\), and its degree-\((2n+1)\) component is
\[
 \frac{q}{(q-q^{-1})^2[2]_q}\,
 \comm{E_\delta}{
      q^{n-1}(q-q^{-1})^{-2n+2}
      E_{(n-1)\delta+\alpha_0}} .
\]
Using \eqref{eq:Damiani-real}, this equals
\[
        q^n(q-q^{-1})^{-2n}E_{n\delta+\alpha_0}.
\]
It is non-zero.  Therefore \(B_{n\delta+\alpha_0}\) has filtration degree
\(2n+1\), and \eqref{eq:root-leading-real0} follows by induction.

The proof of \eqref{eq:root-leading-real1} is identical, apart from the
sign in \eqref{eq:BK-real1-first} and \eqref{eq:BK-real1}.  That sign changes
\(\comm{E_\delta}{E_{(n-1)\delta+\alpha_1}}\) into
\(\comm{E_{(n-1)\delta+\alpha_1}}{E_\delta}\), exactly as required by
\eqref{eq:Damiani-real}.  Thus
\[
        \init(B_{n\delta+\alpha_1})
        =
        q^n(q-q^{-1})^{-2n}E_{n\delta+\alpha_1}
        \qquad(n\in\N).
\]

It remains to treat the imaginary root vectors.  Let \(n\geq1\).  From the
real-root calculation just proved, the first two terms in
\eqref{eq:BK-imaginary} have filtration degree at most \(2n\).  More
precisely, their degree-\(2n\) component is
\[
 q^{n-1}(q-q^{-1})^{-2n+2}
 \bigl(q^{-2}E_{(n-1)\delta+\alpha_1}A
       -AE_{(n-1)\delta+\alpha_1}\bigr).
\]
By \eqref{eq:Damiani-imaginary}, this is
\[
        q^{n-1}(q-q^{-1})^{2-2n}E_{n\delta}.
\]

Every summand in the final sum of \eqref{eq:BK-imaginary} has filtration
degree at most
\[
        (2\ell+1)+\bigl(2(n-\ell-2)+1\bigr)=2n-2,
\]
and therefore cannot contribute to the degree-\(2n\) component.  Since
\(E_{n\delta}\neq0\), we obtain \eqref{eq:root-leading-imaginary}, and
\(B_{n\delta}\) has filtration degree \(2n\).  The proposition follows.
\end{proof}

\begin{remark}
Proposition~\ref{prop:root-leading-U} is the analogue, in the present
normalisation and at the fixed non-root parameter \(q\), of the filtered
comparison in \cite[Proposition~4.4]{BaseilhacKolb2020}.  We have included
the proof because the scalar factors are needed later.
\end{remark}

Combining Proposition~\ref{prop:root-leading-U} with the \(q\)-shuffle
formulae \eqref{eq:Damiani-shuffle-0}--\eqref{eq:Damiani-shuffle-delta}
gives the following explicit initial forms in the \(q\)-shuffle algebra.

\begin{corollary}
\label{cor:root-leading-shuffle}
For \(n\in\N\),
\begin{align}
\init(B_{n\delta+\alpha_0})
 &=q^{-n}xC_n,
\label{eq:BK-leading-xC}\\
\init(B_{n\delta+\alpha_1})
 &=q^{-n}C_ny.
\label{eq:BK-leading-Cy}
\end{align}
For \(n\geq1\),
\begin{equation}
\label{eq:BK-leading-C}
        \init(B_{n\delta})
        =-q^{-n-1}(q-q^{-1})C_n .
\end{equation}
If, only as a scalar convention, one sets
\[
        B_{0\delta}=q^{-2}-1=-q^{-1}(q-q^{-1}),
\]
then \eqref{eq:BK-leading-C} is also valid for \(n=0\), with \(C_0=1\).
\end{corollary}

\begin{proof}
For the \(\alpha_0\)-family, equations
\eqref{eq:root-leading-real0} and \eqref{eq:Damiani-shuffle-0} give
\[
 q^n(q-q^{-1})^{-2n}
 q^{-2n}(q-q^{-1})^{2n}xC_n
 =
 q^{-n}xC_n.
\]
This proves \eqref{eq:BK-leading-xC}.  The proof of
\eqref{eq:BK-leading-Cy} is the same, using
\eqref{eq:root-leading-real1} and \eqref{eq:Damiani-shuffle-1}.  For the
imaginary root vectors, equations \eqref{eq:root-leading-imaginary} and
\eqref{eq:Damiani-shuffle-delta} give
\[
 q^{n-1}(q-q^{-1})^{2-2n}
 \bigl(-q^{-2n}(q-q^{-1})^{2n-1}C_n\bigr)
 =
 -q^{-n-1}(q-q^{-1})C_n.
\]
This is \eqref{eq:BK-leading-C}.  The final assertion follows by substituting
\(n=0\) and \(C_0=1\).
\end{proof}

For a positive root \(\beta\in\Phi_+\), write
\[
 B_\beta=
 \begin{cases}
 B_{n\delta+\alpha_0}, & \beta=n\delta+\alpha_0,\quad n\in\N,\\[1mm]
 B_{n\delta+\alpha_1}, & \beta=n\delta+\alpha_1,\quad n\in\N,\\[1mm]
 B_{n\delta}, & \beta=n\delta,\quad n\geq1.
 \end{cases}
\]
Thus \(\{B_\beta\}_{\beta\in\Phi_+}\) is the full family of
Baseilhac--Kolb root vectors in \(\Oq\).

\begin{theorem}
\label{thm:BK-nonroot}
Let \(\prec\) be any total order on the set \(\Phi_+\) of positive roots of
\(\widehat{\mathfrak{sl}}_2\).  For every finitely supported function
\[
        e:\Phi_+\longrightarrow\N ,
\]
let
\[
        B_\prec^e
        =
        B_{\beta_1}^{e(\beta_1)}
        B_{\beta_2}^{e(\beta_2)}
        \cdots
        B_{\beta_r}^{e(\beta_r)},
\]
where
\[
        \beta_1\prec\beta_2\prec\cdots\prec\beta_r
\]
are the roots in the support of \(e\).  If \(e=0\), set
\(B_\prec^e=1\).  Then the set
\[
        \{\,B_\prec^e\mid e:\Phi_+\to\N
        \text{ finitely supported}\,\}
\]
is an \(\F\)-basis of \(\Oq\).
\end{theorem}

Equivalently, the Baseilhac--Kolb root vectors give a PBW basis of
\(\Oq\) in every total order on \(\Phi_+\), under the sole assumption that
\(q\) is not a root of unity.  In particular, this proves
\cite[Conjecture~16.1]{Terwilliger2022OqACE}.

\begin{proof}
We apply the filtered PBW lifting lemma to the word-length filtration of
\(\Oq\).  By Proposition~\ref{prop:root-leading-U}, every
Baseilhac--Kolb root vector has positive filtration degree and has initial
form equal to a non-zero scalar multiple of the corresponding Damiani root
vector:
\[
        \init(B_\beta)=a_\beta E_\beta
        \qquad(\beta\in\Phi_+),
\]
where \(a_\beta\in\F^\times\).  More explicitly,
\[
 a_{n\delta+\alpha_0}
 =
 a_{n\delta+\alpha_1}
 =
 q^n(q-q^{-1})^{-2n}
 \qquad(n\in\N),
\]
and
\[
        a_{n\delta}
        =
        q^{\,n-1}(q-q^{-1})^{2-2n}
        \qquad(n\geq1).
\]
These scalars are non-zero by Lemma~\ref{lem:nonvanishing-scalars}.

Through the isomorphism
\[
        \gr\Oq\simeq\Uplus
\]
from Theorem~\ref{thm:gr-isomorphism}, the elements \(E_\beta\) are exactly
Damiani's root vectors in the associated graded algebra.  By
Theorem~\ref{thm:any-Damiani}, for the chosen total order \(\prec\), the
ordered Damiani monomials
\[
        E_{\beta_1}^{e(\beta_1)}
        E_{\beta_2}^{e(\beta_2)}
        \cdots
        E_{\beta_r}^{e(\beta_r)}
        \qquad
        (\beta_1\prec\cdots\prec\beta_r)
\]
form a basis of \(\Uplus\simeq\gr\Oq\).

Therefore the hypotheses of Lemma~\ref{lem:filtered-lifting} are satisfied
with the index set \(\Lambda=\Phi_+\), the lifting elements \(u_\beta=B_\beta\),
and the graded elements \(v_\beta=E_\beta\).  The lemma lifts the ordered
Damiani PBW basis of \(\gr\Oq\) to the ordered Baseilhac--Kolb monomials in
\(\Oq\).  Hence the monomials \(B_\prec^e\) form an \(\F\)-basis of
\(\Oq\), as claimed.
\end{proof}





\section{Initial forms of the alternating generators}
\label{sec:alternating-initial}

We now compute the initial forms of the alternating generators of \(\Oq\).
Throughout this section, initial forms are identified with elements of
\(\Uplus\), and then with their images in the \(q\)-shuffle algebra.  Thus,
when an initial form is written as a word expression in \(x,y\), products are
computed using the \(q\)-shuffle product \(\star\).

For $u\in F_d\Oq$, write
\[
        [u]_d=u+F_{d-1}\Oq\in F_d\Oq/F_{d-1}\Oq.
\]
This class is allowed to be zero.  We use the notation $\init(u)$ only after
showing that $[u]_d\neq0$ for the asserted value of $d$.

Recall the canonical central reduction
\[
        \gamma:\Aq\longrightarrow\Oq .
\]
The alternating generators of \(\Oq\) are
\[
 W_{-n}=\gamma(\cW_{-n}),\qquad
 W_{n+1}=\gamma(\cW_{n+1}),\qquad
 G_{n+1}=\gamma(\cG_{n+1}),\qquad
 \widetilde G_{n+1}=\gamma(\ctG_{n+1})
 \qquad(n\in\N).
\]
This is Terwilliger's notation from
\cite[Definition~11.5]{Terwilliger2022OqACE}.  We also use the scalar
convention
\begin{equation}
\label{eq:G0-Oq}
        G_0=\widetilde G_0=-(q-q^{-1})[2]_q^2 .
\end{equation}
The elements \(G_0\) and \(\widetilde G_0\) are scalars, not additional
generators.

We shall use two relations obtained from the defining relations of \(\Aq\)
after applying \(\gamma\).  For \(n\geq1\),
\begin{align}
\qcomm{\widetilde G_n}{W_0}
 &=\rho_0W_{-n}-\rho_0W_n,
\label{eq:Oq-tG-W0}\\
\qcomm{W_1}{\widetilde G_n}
 &=\rho_0W_{n+1}-\rho_0W_{-(n-1)},
\label{eq:Oq-W1-tG}
\end{align}
where \(\rho_0=-(q^2-q^{-2})^2\).  These are the specialisations of
\eqref{eq:Aq2} and \eqref{eq:Aq3}.  We shall also use the
antiautomorphism \(\dagger\) of \(\Oq\), which fixes \(W_0,W_1\) and
interchanges the two imaginary alternating families:
\begin{equation}
\label{eq:dagger-alternating}
        G_n^\dagger=\widetilde G_n,\qquad
        \widetilde G_n^\dagger=G_n
        \qquad(n\geq1).
\end{equation}
See \cite[Theorem~11.6 and Lemma~11.8]{Terwilliger2022OqACE}.  The same
statements are harmless for \(n=0\), since \(G_0=\widetilde G_0\) is a
scalar.

Finally, we use the scalar convention introduced after
Corollary~\ref{cor:root-leading-shuffle}:
\begin{equation}
\label{eq:B0delta-convention-alt}
        B_{0\delta}=q^{-2}-1=-q^{-1}(q-q^{-1}).
\end{equation}
With this convention, the formula
\[
        \init(B_{m\delta})
        =
        -q^{-m-1}(q-q^{-1})C_m
\]
holds for every \(m\in\N\), including \(m=0\).

\subsection{The root-vector recursion and the family \texorpdfstring{\(\widetilde G_n\)}{G-tilde n}}

The main input for the computation is Terwilliger's recursion relating the
imaginary Baseilhac--Kolb root vectors to the family
\(\{\widetilde G_n\}_{n\geq0}\).

\begin{lemma}[{{\cite[Equation~(133)]{Terwilliger2022OqACE}}; see also
{\cite[Lemma~12.5]{Terwilliger2022OqACE}}}]
\label{lem:root-tG-recursion}
For \(n\geq1\),
\begin{align}
0={}&[n]_qB_{n\delta}\widetilde G_0+\sum_{\substack{j+k+2\ell+1=n\\j,k,\ell\geq0}}
(-1)^\ell\binom{k+\ell}{\ell}
[2n-j]_q[2]_q^{k+1}B_{j\delta}\widetilde G_{k+1}.
\label{eq:root-tG-recursion}
\end{align}
\end{lemma}



\begin{proposition}
\label{prop:tG-leading}
For every $n\in\N$, one has $\widetilde G_n\in F_{2n}\Oq$ and
\begin{equation}
\label{eq:tG-leading}
\init(\widetilde G_n)
 =
 (-1)^{n+1}(q-q^{-1})q^{-n}[2]_q^{2-n}
 \widetilde G_n^{\sh}.
\end{equation}
Consequently,
\[
        \deg_F\widetilde G_n=2n .
\]
\end{proposition}

\begin{proof}
We argue by induction on $n$.  For $n=0$, the assertion is the scalar
identity \eqref{eq:G0-Oq}, since $\widetilde G_0^{\sh}=1$.

Let $n\geq1$, and assume the assertion for all smaller indices.  In
\eqref{eq:root-tG-recursion}, the only occurrence of $\widetilde G_n$ is
the term with $(j,k,\ell)=(0,n-1,0)$, whose coefficient is
\[
        [2n]_q[2]_q^nB_{0\delta}.
\]
This scalar is non-zero.  Every other alternating generator in the recursion
has index smaller than $n$.  By the induction hypothesis and
Corollary~\ref{cor:root-leading-shuffle}, the term indexed by
$j,k,\ell$ belongs to filtration degree at most
\[
        2j+2(k+1)=2n-4\ell.
\]
Solving the recursion for $\widetilde G_n$ therefore gives
$\widetilde G_n\in F_{2n}\Oq$.

We now pass to $F_{2n}\Oq/F_{2n-1}\Oq$.  Put $r=k+1$.  The terms with
$\ell>0$ disappear, and we obtain
\begin{equation}
\label{eq:top-recursion-Oq}
0=
\sum_{r=0}^{n}
[n+r]_q[2]_q^r\,
[ B_{(n-r)\delta}]_{2(n-r)}\star
[\widetilde G_r]_{2r}.
\end{equation}
For $r<n$, the second factor is known by induction; the first factor is
known from Corollary~\ref{cor:root-leading-shuffle}, including the scalar
case $n-r=0$.

Set
\[
        X_n=
        (-1)^{n+1}(q-q^{-1})q^{-n}[2]_q^{2-n}
        \widetilde G_n^{\sh}.
\]
Substituting the known classes for $r<n$ and substituting $X_n$ for the
unknown class $[\widetilde G_n]_{2n}$, the coefficient of
$C_{n-r}\star\widetilde G_r^{\sh}$ becomes, for every $0\leq r\leq n$,
\[
        (-1)^r q^{-n-1}(q-q^{-1})^2[2]_q^2[n+r]_q.
\]
After removing the common non-zero factor, equation
\eqref{eq:top-recursion-Oq} is exactly
\[
        \sum_{r=0}^{n}(-1)^r[n+r]_q
        C_{n-r}\star\widetilde G_r^{\sh}=0,
\]
which is Theorem~\ref{thm:Catalan-alternating} after the substitution
$i=n-r$ and multiplication by $(-1)^n$.

The coefficient of the unknown class in
\eqref{eq:top-recursion-Oq} is the non-zero scalar
$[2n]_q[2]_q^nB_{0\delta}$.  Hence that equation has a unique solution for
$[\widetilde G_n]_{2n}$, and the preceding calculation shows that
\[
        [\widetilde G_n]_{2n}=X_n.
\]
The right-hand side is non-zero.  Therefore
$\widetilde G_n\notin F_{2n-1}\Oq$, so its filtration degree is $2n$ and
its initial form is $X_n$.
\end{proof}

\begin{remark}[Low-degree check]
For $n=1$, the recursion in Lemma~\ref{lem:root-tG-recursion} gives the
exact identity
\begin{equation}
\label{eq:tG1-Bdelta}
        \widetilde G_1=-qB_\delta=[W_0,W_1]_q.
\end{equation}
Applying $\dagger$ gives
\begin{equation}
\label{eq:G1-low-check}
        G_1=[W_1,W_0]_q.
\end{equation}
\end{remark}

\subsection{The remaining alternating families}

\begin{proposition}
\label{prop:W-leading}
For every $n\in\N$,
\begin{align}
\init(W_{-n})
 &=(-1)^nq^{-n}[2]_q^{-n}W_{-n}^{\sh},
\label{eq:Wminus-leading}\\
\init(W_{n+1})
 &=(-1)^nq^{-n}[2]_q^{-n}W_{n+1}^{\sh}.
\label{eq:Wplus-leading}
\end{align}
Consequently,
\[
        \deg_F W_{-n}=\deg_F W_{n+1}=2n+1 .
\]
\end{proposition}

\begin{proof}
We prove both assertions simultaneously by induction on $n$.  For $n=0$,
\[
        \init(W_0)=x=W_0^{\sh},\qquad
        \init(W_1)=y=W_1^{\sh}.
\]

Let $n\geq1$ and assume the assertions for smaller indices.  By
Proposition~\ref{prop:tG-leading}, $\widetilde G_n\in F_{2n}\Oq$.
Equations \eqref{eq:Oq-tG-W0} and \eqref{eq:Oq-W1-tG}, together with the
induction hypothesis for $W_n$ and $W_{-(n-1)}$, first show that
\[
        W_{-n},W_{n+1}\in F_{2n+1}\Oq.
\]
Passing to the degree-$(2n+1)$ quotient gives
\begin{align}
\qcomm{\init(\widetilde G_n)}{x}
 &=\rho_0[W_{-n}]_{2n+1},
\label{eq:top-Wminus}\\
\qcomm{y}{\init(\widetilde G_n)}
 &=\rho_0[W_{n+1}]_{2n+1}.
\label{eq:top-Wplus}
\end{align}
The lower-index $W$-terms have disappeared because they have degree
$2n-1$.

Substituting Proposition~\ref{prop:tG-leading} and using
\eqref{eq:qshuffle-Wminus}--\eqref{eq:qshuffle-Wplus}, each left-hand side
is respectively
\[
 (-1)^{n+1}(q-q^{-1})^2q^{-n}[2]_q^{2-n}W_{-n}^{\sh},
 \qquad
 (-1)^{n+1}(q-q^{-1})^2q^{-n}[2]_q^{2-n}W_{n+1}^{\sh}.
\]
Since
\[
        \rho_0=-(q-q^{-1})^2[2]_q^2,
\]
division by $\rho_0$ yields
\[
 [W_{-n}]_{2n+1}=(-1)^nq^{-n}[2]_q^{-n}W_{-n}^{\sh},
\]
\[
 [W_{n+1}]_{2n+1}=(-1)^nq^{-n}[2]_q^{-n}W_{n+1}^{\sh}.
\]
Both classes are non-zero.  Thus the two elements have exact filtration
degree $2n+1$, and the displayed classes are their initial forms.
\end{proof}

\begin{proposition}
\label{prop:G-leading}
For every \(n\in\N\),
\begin{equation}
\label{eq:G-leading}
\init(G_n)
 =
 (-1)^{n+1}(q-q^{-1})q^{-n}[2]_q^{2-n}G_n^{\sh}.
\end{equation}
Consequently,
\[
        \deg_F G_n=2n .
\]
\end{proposition}

\begin{proof}
For \(n=0\), this is the scalar identity
\(G_0=-(q-q^{-1})[2]_q^2\), since \(G_0^{\sh}=1\).

Let \(n\geq1\).  The antiautomorphism \(\dagger\) fixes \(W_0\) and
\(W_1\), and therefore preserves the word-length filtration.  Its induced
antiautomorphism on \(\gr\Oq\) fixes \(\overline{W_0}\) and
\(\overline{W_1}\).  Under the \(q\)-shuffle realisation, this induced map
is the word-reversal antiautomorphism.  By
\eqref{eq:dagger-alternating}, \(G_n=\widetilde G_n^\dagger\), while word
reversal sends \(\widetilde G_n^{\sh}\) to \(G_n^{\sh}\).  Applying
\(\dagger\) to the initial-form formula in
Proposition~\ref{prop:tG-leading} gives \eqref{eq:G-leading}.  The scalar
is non-zero, so \(\deg_FG_n=2n\).
\end{proof}

Combining the three preceding propositions gives the main result of this
section.

\begin{theorem}
\label{thm:all-leading}
For every \(n\in\N\),
\begin{align}
\init(W_{-n})
 &=(-1)^nq^{-n}[2]_q^{-n}W_{-n}^{\sh},
\label{eq:all-leading-Wminus}\\
\init(W_{n+1})
 &=(-1)^nq^{-n}[2]_q^{-n}W_{n+1}^{\sh},
\label{eq:all-leading-Wplus}\\
\init(G_n)
 &=(-1)^{n+1}(q-q^{-1})q^{-n}[2]_q^{2-n}G_n^{\sh},
\label{eq:all-leading-G}\\
\init(\widetilde G_n)
 &=(-1)^{n+1}(q-q^{-1})q^{-n}[2]_q^{2-n}
   \widetilde G_n^{\sh}.
\label{eq:all-leading-tG}
\end{align}
In particular,
\[
        \deg_F W_{-n}=\deg_F W_{n+1}=2n+1,
        \qquad
        \deg_F G_n=\deg_F\widetilde G_n=2n .
\]
\end{theorem}

\begin{proof}
Equations \eqref{eq:all-leading-Wminus} and \eqref{eq:all-leading-Wplus}
are Proposition~\ref{prop:W-leading}.  Equation
\eqref{eq:all-leading-G} is Proposition~\ref{prop:G-leading}, and
\eqref{eq:all-leading-tG} is Proposition~\ref{prop:tG-leading}.  The degree
statements are part of the same propositions.
\end{proof}







\section{Twelve alternating PBW bases}
\label{sec:twelve}

We now lift Terwilliger's alternating PBW bases of \(\Uplus\) to the
\(q\)-Onsager algebra \(\Oq\).  We use the following block notation:
\[
 W^-=\{W_{-i}\mid i\in\N\},\qquad
 W^+=\{W_{k+1}\mid k\in\N\},
\]
and
\[
 G=\{G_{j+1}\mid j\in\N\},\qquad
 \widetilde G=\{\widetilde G_{j+1}\mid j\in\N\}.
\]
Thus \(G\) and \(\widetilde G\) do not include the scalar elements
\(G_0\) and \(\widetilde G_0\).

A linear order on one of the three-family unions
\[
        W^-\cup \widetilde G\cup W^+,
        \qquad
        W^-\cup G\cup W^+
\]
is said to satisfy a displayed block order if every element of an earlier
block is smaller than every element of a later block.  No restriction is
imposed on the order inside a single block.

\begin{theorem}[Alternating PBW bases of \(\Oq\)]
\label{thm:twelve-PBW}
For \(H=\widetilde G\) and for \(H=G\), the three families
\[
        W^-,\qquad H,\qquad W^+
\]
give PBW bases of \(\Oq\) in each of the six block orders
\begin{equation}
\label{eq:six-Oq-orders}
\begin{array}{lll}
W^-<H<W^+,&
W^+<H<W^-,&
W^+<W^-<H,\\[1mm]
W^-<W^+<H,&
H<W^+<W^-,&
H<W^-<W^+.
\end{array}
\end{equation}
Equivalently, after choosing any linear order inside each of the three
blocks, the ordered monomials in the resulting ordered set form an
\(\F\)-basis of \(\Oq\).
\end{theorem}

\begin{proof}
We give the proof for \(H=\widetilde G\); the proof for \(H=G\) is
identical.

Fix one of the six block orders in \eqref{eq:six-Oq-orders}, and fix an
arbitrary linear order inside each block.  By
Theorem~\ref{thm:all-leading}, the initial forms of the generators in the
three families are non-zero scalar multiples of the corresponding
alternating words:
\[
 \init(W_{-i})\in\F^\times W_{-i}^{\sh},\qquad
 \init(\widetilde G_{j+1})\in\F^\times\widetilde G_{j+1}^{\sh},\qquad
 \init(W_{k+1})\in\F^\times W_{k+1}^{\sh}.
\]
The scalars are non-zero by Lemma~\ref{lem:nonvanishing-scalars}.  Through
the isomorphism \(\gr\Oq\simeq\Uplus\) from
Theorem~\ref{thm:gr-isomorphism}, these initial forms are elements of
\(\Uplus\).

By Theorem~\ref{thm:qshuffle-twelve}, the ordered monomials in
\[
 \{W_{-i}^{\sh}\}_{i\in\N},\qquad
 \{\widetilde G_{j+1}^{\sh}\}_{j\in\N},\qquad
 \{W_{k+1}^{\sh}\}_{k\in\N}
\]
form a basis of \(\Uplus\) for the chosen block order and the chosen
internal orders.  Therefore the hypotheses of the filtered PBW lifting
lemma, Lemma~\ref{lem:filtered-lifting}, are satisfied.  The corresponding
ordered monomials in
\[
 \{W_{-i}\}_{i\in\N},\qquad
 \{\widetilde G_{j+1}\}_{j\in\N},\qquad
 \{W_{k+1}\}_{k\in\N}
\]
form a basis of \(\Oq\).

Replacing \(\widetilde G_{j+1}\) by \(G_{j+1}\) and using the \(G\)-part of
Theorem~\ref{thm:qshuffle-twelve}, together with
Theorem~\ref{thm:all-leading}, gives the same conclusion for \(H=G\).
\end{proof}

Terwilliger's Conjecture~16.2 in
\cite{Terwilliger2022OqACE}  is the case \(H=\widetilde G\) of
Theorem~\ref{thm:twelve-PBW}.

\begin{remark}
Theorem~\ref{thm:twelve-PBW} gives twelve PBW bases: six block orders with
the middle imaginary family \(\widetilde G\), and six block orders with the
middle imaginary family \(G\).  The order inside a single block is
irrelevant for existence of the basis.  Indeed, the corresponding
\(q\)-shuffle theorem already allows arbitrary internal orders, and the
members of each alternating family in \(\Oq\) commute pairwise by
\cite[Theorem~11.6]{Terwilliger2022OqACE}.
\end{remark}

\begin{remark}
All twelve PBW bases have the same filtered Hilbert series.  In each basis
there are two generators of filtration degree \(2m-1\), namely
\(W_{-(m-1)}\) and \(W_m\), and one generator of filtration degree \(2m\),
namely either \(G_m\) or \(\widetilde G_m\).  Hence the filtered Hilbert
series is
\[
        \prod_{m\geq1}\frac{1}{(1-t^{2m})}
        \prod_{m\geq1}\frac{1}{(1-t^{2m-1})^2}.
\]
This agrees with the Hilbert series of \(\gr\Oq\simeq\Uplus\).
\end{remark}



\section{Central specialisations of the alternating central extension}
\label{sec:central-specialisations}

We next show that the alternating PBW bases of
Theorem~\ref{thm:twelve-PBW} are stable under arbitrary central
specialisation of the alternating central extension \(\Aq\).

Terwilliger proved that the standard tensor product factorisation of
\(\Aq\) may be written in the following explicit form.  Put \(z_0=1\).
There is an algebra isomorphism
\[
        \Phi:\Aq\longrightarrow \Oq\otimes\F[z_1,z_2,\ldots]
\]
such that, for \(n\in\N\),
\begin{align}
\Phi(\cW_{-n})
 &=\sum_{k=0}^{n}W_{k-n}\otimes z_k,
\label{eq:central-convolution-Wminus}\\
\Phi(\cW_{n+1})
 &=\sum_{k=0}^{n}W_{n+1-k}\otimes z_k,
\label{eq:central-convolution-Wplus}\\
\Phi(\cG_n)
 &=\sum_{k=0}^{n}G_{n-k}\otimes z_k,
\label{eq:central-convolution-G}\\
\Phi(\ctG_n)
 &=\sum_{k=0}^{n}\widetilde G_{n-k}\otimes z_k.
\label{eq:central-convolution-tG}
\end{align}
Here \(W_{k-n}=W_{-(n-k)}\).  The index-zero terms
\[
        G_0=\widetilde G_0=-(q-q^{-1})[2]_q^2
\]
are scalars.  The isomorphism \(\Phi\), and the convolution formulae above,
are from \cite[Theorem~13.5 and Lemma~13.1]
{Terwilliger2022OqACE}.

\subsection{Central coordinates and central fibres}

Several systems of central coordinates occur in the literature.  We record
the precise conventions used below.

First, let \(\XiC_n\) denote the central elements denoted $Z_n$ in
\cite[Definition~8.6]{Terwilliger2021ACE}; we rename them here to avoid
confusion with the convolution coordinates below.  With \(h=[2]_q\),
the generating function in that reference is
\[
        Z(t)=\sum_{n\geq0}\XiC_nh^nt^n.
\]
Under the notation of \cite[Definition~8.1 and Note~8.2]
{Terwilliger2022OqACE}, this same generating function is denoted
\(\Psi(t)\).
The normalised central elements of
\cite[Definitions~8.3 and~8.4]{Terwilliger2022OqACE} are
\begin{equation}
\label{eq:Xi-Zvee}
        \Zvee_n=h^{n-2}\XiC_n
        \qquad(n\geq0).
\end{equation}
In particular, \(\Zvee_0=1\).

Secondly, Baseilhac and Belliard use central elements denoted
\(\Delta_n\).  When \(\operatorname{char}\F\neq2\), Terwilliger's comparison
formula gives
\begin{equation}
\label{eq:Delta-Xi}
 \Delta_n
 =
 -2\frac{q-q^{-1}}{q^n+q^{-n}}\XiC_n
 =
 -2\frac{(q-q^{-1})h^{2-n}}{q^n+q^{-n}}\Zvee_n
 \qquad(n\geq1).
\end{equation}
See \cite[Remark~8.17]{Terwilliger2021ACE}.  Under the additional
assumption \(\operatorname{char}\F\neq2\), every scalar factor in
\eqref{eq:Delta-Xi} is non-zero by
Lemma~\ref{lem:nonvanishing-scalars}.  In characteristic \(2\), the
coordinates \(\Delta_n\) do not provide an equivalent system of central
coordinates; in that case we use the coordinates \(\Zvee_n\) or the
convolution coordinates below.

Finally, define the convolution central elements \(\Zconv_n\) by
\begin{equation}
\label{eq:convolution-central-coordinate}
        \Phi(\Zconv_n)=1\otimes z_n,
        \qquad
        \Zconv_0=1 .
\end{equation}
The elements \(\Zconv_n\) are central in \(\Aq\).  Terwilliger proved that
the two systems
\[
        \{\Zvee_n\}_{n\geq1},
        \qquad
        \{\Zconv_n\}_{n\geq1}
\]
are related by a triangular polynomial change of variables with
invertible diagonal coefficients.  More precisely, for \(n\geq1\),
\begin{equation}
\label{eq:Zvee-Zconv-triangular}
 \Zvee_n
 =
 h^n(q^n+q^{-n})\Zconv_n
 +P_n(\Zconv_1,\ldots,\Zconv_{n-1}),
\end{equation}
where \(P_n\) has zero constant term and is homogeneous of weighted degree
\(n\) when \(\deg\Zconv_k=k\).  See
\cite[Lemma~13.10(i)]{Terwilliger2022OqACE}.  Since
\(h^n(q^n+q^{-n})\neq0\), this triangular system can be inverted
recursively.  Thus \(\Zconv_n\) is a polynomial in
\(\Zvee_1,\ldots,\Zvee_n\), with leading coefficient
\(h^{-n}(q^n+q^{-n})^{-1}\).

Consequently, prescribing the values of the \(\Zvee_n\) is equivalent to
prescribing the values of the \(\Zconv_n\).  If
\(\operatorname{char}\F\neq2\), prescribing the values of the
\(\Delta_n\) is equivalent to prescribing either of these two systems.  The
three systems generate the same centre in characteristic different from
\(2\), but they are not termwise identical.

Let
\[
        \boldsymbol\zeta=(\zeta_1,\zeta_2,\ldots)\in\F^{\N_{\geq1}},
        \qquad
        \zeta_0=1 .
\]
Evaluation at \(z_k=\zeta_k\) defines an algebra homomorphism
\[
        \operatorname{ev}_{\boldsymbol\zeta}:
        \F[z_1,z_2,\ldots]\longrightarrow\F .
\]
Its kernel is the ideal generated by the elements $z_k-\zeta_k$
($k\geq1$).  Equivalently, we specialise the central elements by
\[
        \Zconv_k\longmapsto \zeta_k
        \qquad(k\geq1).
\]
The corresponding central fibre is
\[
        \Aqfib{\boldsymbol\zeta}
        =
        \Aq\big/
        \left\langle
        \Zconv_k-\zeta_k\mid k\geq1
        \right\rangle .
\]
Using \(\Phi\), we identify this fibre with \(\Oq\) through the composite
map
\[
        \Aq
        \xrightarrow{\ \Phi\ }
        \Oq\otimes\F[z_1,z_2,\ldots]
        \xrightarrow{\ \operatorname{id}\otimes
        \operatorname{ev}_{\boldsymbol\zeta}\ }
        \Oq .
\]
Under this identification, the images of the alternating generators of
\(\Aq\) are the following elements of \(\Oq\):
\begin{align}
W_{-n}^{(\boldsymbol\zeta)}
 &=W_{-n}+\sum_{k=1}^{n}\zeta_kW_{-(n-k)},
\label{eq:special-Wminus}\\
W_{n+1}^{(\boldsymbol\zeta)}
 &=W_{n+1}+\sum_{k=1}^{n}\zeta_kW_{n+1-k},
\label{eq:special-Wplus}\\
G_n^{(\boldsymbol\zeta)}
 &=G_n+\sum_{k=1}^{n}\zeta_kG_{n-k},
\label{eq:special-G}\\
\widetilde G_n^{(\boldsymbol\zeta)}
 &=\widetilde G_n+\sum_{k=1}^{n}\zeta_k\widetilde G_{n-k}.
\label{eq:special-tG}
\end{align}
For \(G_n^{(\boldsymbol\zeta)}\) and
\(\widetilde G_n^{(\boldsymbol\zeta)}\), the term with \(k=n\) involves the
scalar \(G_0=\widetilde G_0=-(q-q^{-1})[2]_q^2\).

\begin{proposition}
\label{prop:special-leading}
For every $\F$-valued central character \(\boldsymbol\zeta\), the specialised
alternating generators have the same initial forms as the canonical
alternating generators.  Namely, for \(n\in\N\),
\begin{align}
\init(W_{-n}^{(\boldsymbol\zeta)})
 &=\init(W_{-n}),
\label{eq:special-leading-Wminus}\\
\init(W_{n+1}^{(\boldsymbol\zeta)})
 &=\init(W_{n+1}),
\label{eq:special-leading-Wplus}\\
\init(G_n^{(\boldsymbol\zeta)})
 &=\init(G_n),
\label{eq:special-leading-G}\\
\init(\widetilde G_n^{(\boldsymbol\zeta)})
 &=\init(\widetilde G_n).
\label{eq:special-leading-tG}
\end{align}
\end{proposition}

\begin{proof}
For \(n=0\), there are no correction terms in
\eqref{eq:special-Wminus}--\eqref{eq:special-tG}, so the assertion is
immediate.

Let \(n\geq1\).  By Theorem~\ref{thm:all-leading}, the leading terms of
\(W_{-n}\) and \(W_{n+1}\) have filtration degree \(2n+1\).  For
\(1\leq k\leq n\),
\[
        \deg_F W_{-(n-k)}=2(n-k)+1<2n+1,
\]
and
\[
        \deg_F W_{n+1-k}=2(n-k)+1<2n+1 .
\]
Thus every correction term in
\eqref{eq:special-Wminus} and \eqref{eq:special-Wplus} has strictly smaller
filtration degree than the leading summand.

Similarly, the leading terms of \(G_n\) and \(\widetilde G_n\) have
filtration degree \(2n\).  For \(1\leq k\leq n\),
\[
        \deg_F G_{n-k}=\deg_F\widetilde G_{n-k}=2(n-k)<2n .
\]
Hence every correction term in \eqref{eq:special-G} and
\eqref{eq:special-tG} has strictly smaller filtration degree.  The initial
forms are therefore unchanged.
\end{proof}

For the specialised fibre, put
\[
 W^-_{\boldsymbol\zeta}
 =
 \{W_{-i}^{(\boldsymbol\zeta)}\mid i\in\N\},
 \qquad
 W^+_{\boldsymbol\zeta}
 =
 \{W_{k+1}^{(\boldsymbol\zeta)}\mid k\in\N\},
\]
and
\[
 G_{\boldsymbol\zeta}
 =
 \{G_{j+1}^{(\boldsymbol\zeta)}\mid j\in\N\},
 \qquad
 \widetilde G_{\boldsymbol\zeta}
 =
 \{\widetilde G_{j+1}^{(\boldsymbol\zeta)}\mid j\in\N\}.
\]

\begin{theorem}[PBW bases in every $\F$-valued central fibre]
\label{thm:central-fibre-PBW}
For every $\F$-valued central character \(\boldsymbol\zeta\), and for
\(H_{\boldsymbol\zeta}=\widetilde G_{\boldsymbol\zeta}\) or
\(H_{\boldsymbol\zeta}=G_{\boldsymbol\zeta}\), the three specialised
families
\[
        W^-_{\boldsymbol\zeta},
        \qquad
        H_{\boldsymbol\zeta},
        \qquad
        W^+_{\boldsymbol\zeta}
\]
give PBW bases of the central fibre \(\Aqfib{\boldsymbol\zeta}\simeq\Oq\) in
each of the six block orders
\[
\begin{array}{lll}
W^-_{\boldsymbol\zeta}<H_{\boldsymbol\zeta}<W^+_{\boldsymbol\zeta},&
W^+_{\boldsymbol\zeta}<H_{\boldsymbol\zeta}<W^-_{\boldsymbol\zeta},&
W^+_{\boldsymbol\zeta}<W^-_{\boldsymbol\zeta}<H_{\boldsymbol\zeta},\\[1mm]
W^-_{\boldsymbol\zeta}<W^+_{\boldsymbol\zeta}<H_{\boldsymbol\zeta},&
H_{\boldsymbol\zeta}<W^+_{\boldsymbol\zeta}<W^-_{\boldsymbol\zeta},&
H_{\boldsymbol\zeta}<W^-_{\boldsymbol\zeta}<W^+_{\boldsymbol\zeta}.
\end{array}
\]
As before, no condition is imposed on the order inside a single block.
\end{theorem}

\begin{proof}
We prove the statement for
\(H_{\boldsymbol\zeta}=\widetilde G_{\boldsymbol\zeta}\); the proof for
\(G_{\boldsymbol\zeta}\) is identical.

By Proposition~\ref{prop:special-leading}, the specialised generators have
exactly the same initial forms as the canonical alternating generators.
By Theorem~\ref{thm:all-leading}, these initial forms are non-zero scalar
multiples of the corresponding alternating words in \(\Uplus\).  The
ordered monomials in those alternating words form a basis of \(\Uplus\) by
Theorem~\ref{thm:qshuffle-twelve}.  Since
\(\gr\Oq\simeq\Uplus\), Lemma~\ref{lem:filtered-lifting} applies and lifts
the graded PBW basis to the specialised generators.  Therefore the ordered
monomials in
\[
        W^-_{\boldsymbol\zeta},\qquad
        \widetilde G_{\boldsymbol\zeta},\qquad
        W^+_{\boldsymbol\zeta}
\]
form a basis of the fibre \(\Aqfib{\boldsymbol\zeta}\simeq\Oq\) in each of
the six block orders.  Replacing \(\widetilde G_{\boldsymbol\zeta}\) by
\(G_{\boldsymbol\zeta}\) gives the remaining six bases.
\end{proof}


\section{The conjecture of Baseilhac and Belliard}
\label{sec:BB}

We now apply the central-fibre PBW theorem to the quotients introduced by
Baseilhac and Belliard.  Throughout this section we assume
\[
        \operatorname{char}\F\neq2 .
\]
This assumption is needed in order to use the central coordinates
\(\Delta_n\), since their comparison with Terwilliger's central coordinates
contains the scalar factor \(2\).

\subsection{The normalised parameter}

Let \(\Aq(\rho_0)\) be the current algebra with
\[
        \rho_0=-(q^2-q^{-2})^2 .
\]
This is the normalisation compatible with the \(q\)-Onsager algebra
\(\Oq\) used in the present paper.  Baseilhac and Belliard constructed
central elements
\[
        \Delta_1,\Delta_2,\ldots
\]
in their current algebra.  For a scalar sequence
\[
        \boldsymbol\delta=(\delta_1,\delta_2,\ldots),
\]
they define the quotient
\begin{equation}
\label{eq:BB-quotient-normalised}
\BBqquot{\rho_0}{\boldsymbol\delta}
 =
 \Aq(\rho_0)/
 \left\langle
        \Delta_n-2\delta_n\mid n\geq1
 \right\rangle .
\end{equation}
See \cite[Definition~3.1]{BaseilhacBelliard2017}.

We first identify this quotient with a central fibre of \(\Aq\).  By
\eqref{eq:Delta-Xi}, the relation \(\Delta_n=2\delta_n\) is equivalent to
\begin{equation}
\label{eq:BB-Zvee-values}
 \Zvee_n
 =
 -\frac{(q^n+q^{-n})[2]_q^{\,n-2}}{q-q^{-1}}\,
 \delta_n
 \qquad(n\geq1).
\end{equation}
The coefficient in \eqref{eq:BB-Zvee-values} is non-zero by
Lemma~\ref{lem:nonvanishing-scalars}.  Since the coordinates
\(\{\Zvee_n\}_{n\geq1}\) and \(\{\Zconv_n\}_{n\geq1}\) are related by the
triangular system with invertible diagonal coefficients \eqref{eq:Zvee-Zconv-triangular}, the prescribed values
\eqref{eq:BB-Zvee-values} determine a unique sequence
\[
        \boldsymbol\zeta=(\zeta_1,\zeta_2,\ldots)
\]
such that
\[
        \Zconv_n\longmapsto \zeta_n
        \qquad(n\geq1).
\]
Consequently, \(\BBqquot{\rho_0}{\boldsymbol\delta}\) is the central fibre
\(\Aqfib{\boldsymbol\zeta}\) from
Section~\ref{sec:central-specialisations}.  Under this identification, the
images of the Baseilhac--Belliard generators are exactly the specialised
alternating generators
\[
 W_{-n}^{(\boldsymbol\zeta)},\qquad
 W_{n+1}^{(\boldsymbol\zeta)},\qquad
 G_{n+1}^{(\boldsymbol\zeta)},\qquad
 \widetilde G_{n+1}^{(\boldsymbol\zeta)} .
\]

Let us check the first generator explicitly, since this fixes the
normalisation.  For \(n=1\), equations \eqref{eq:BB-Zvee-values} and
\eqref{eq:Zvee-Zconv-triangular} give
\[
        \Zvee_1=-\frac{\delta_1}{q-q^{-1}},
        \qquad
        \zeta_1=-\frac{\delta_1}{(q-q^{-1})[2]_q^2}.
\]
Using \(G_0=-(q-q^{-1})[2]_q^2\), we obtain
\begin{equation}
\label{eq:BB-low-check}
        G_1^{(\boldsymbol\zeta)}
        =
        G_1+\zeta_1G_0
        =
        G_1+\delta_1 .
\end{equation}
By the low-degree identity \eqref{eq:G1-low-check},
\(G_1=[W_1,W_0]_q\).  Therefore
\[
        G_1^{(\boldsymbol\zeta)}
        =
        [W_1,W_0]_q+\delta_1,
\]
which agrees with \cite[Equation~(3.5)]{BaseilhacBelliard2017}.  Similarly,
\[
        \widetilde G_1^{(\boldsymbol\zeta)}
        =
        [W_0,W_1]_q+\delta_1.
\]

We now state the PBW result in the notation of Baseilhac and Belliard.  The
generators appearing below are the images of the corresponding current
generators in the quotient
\(\BBqquot{\rho_0}{\boldsymbol\delta}\).

\begin{theorem}[Baseilhac--Belliard basis, normalised parameter]
\label{thm:BB-normalised}
Assume that \(\operatorname{char}\F\neq2\).  The monomials
\begin{equation}
\label{eq:BB-monomials}
 W_{-k_1}^{a_1}\cdots W_{-k_N}^{a_N}
 G_{p_1+1}^{b_1}\cdots G_{p_P+1}^{b_P}
 W_{\ell_M+1}^{c_M}\cdots W_{\ell_1+1}^{c_1},
\end{equation}
where
\[
        N,P,M\in\N,
\]
\[
        0\leq k_1<\cdots<k_N,\qquad
        0\leq p_1<\cdots<p_P,\qquad
        0\leq \ell_1<\cdots<\ell_M,
\]
and
\[
        a_i,b_j,c_t\in\Z_{>0},
\]
form an \(\F\)-basis of
\(\BBqquot{\rho_0}{\boldsymbol\delta}\).  Empty products are interpreted
as \(1\).  Equivalently, one may allow non-negative exponents in
\eqref{eq:BB-monomials} after deleting the factors with exponent zero.
\end{theorem}

\begin{proof}
As shown above,
\(\BBqquot{\rho_0}{\boldsymbol\delta}\) is the central fibre
\(\Aqfib{\boldsymbol\zeta}\) associated with the convolution central character
\(\Zconv_n\mapsto\zeta_n\).  Theorem~\ref{thm:central-fibre-PBW}, applied
to the block order
\[
        W^-_{\boldsymbol\zeta}
        <
        G_{\boldsymbol\zeta}
        <
        W^+_{\boldsymbol\zeta},
\]
says that the ordered monomials in the specialised families
\[
        \{W_{-i}^{(\boldsymbol\zeta)}\}_{i\in\N},
        \qquad
        \{G_{j+1}^{(\boldsymbol\zeta)}\}_{j\in\N},
        \qquad
        \{W_{k+1}^{(\boldsymbol\zeta)}\}_{k\in\N}
\]
form a basis.

Choose increasing index order in the \(W^-\)-block and in the \(G\)-block,
and choose decreasing index order in the \(W^+\)-block.  Written in the
notation of the quotient, the resulting ordered monomials are precisely
those in \eqref{eq:BB-monomials}.  Hence the displayed monomials form an
\(\F\)-basis of
\(\BBqquot{\rho_0}{\boldsymbol\delta}\).
\end{proof}

\subsection{Reduction of a general non-zero parameter}

Baseilhac and Belliard formulate their current algebra with an arbitrary
non-zero parameter \(\rho\).  We now reduce this case to the normalised
parameter \(\rho_0\).

Let \(\Aq(\rho)\) denote the Baseilhac--Belliard current algebra with
parameter \(\rho\), and define
\begin{equation}
\label{eq:BB-quotient-rho}
\BBqquot{\rho}{\boldsymbol\delta}
 =
 \Aq(\rho)/
 \left\langle
        \Delta_n^{(\rho)}-2\delta_n\mid n\geq1
 \right\rangle .
\end{equation}
Here \(\Delta_n^{(\rho)}\) denotes the central element \(\Delta_n\) in the
current algebra with parameter \(\rho\).

\begin{lemma}[Scaling of the parameter]
\label{lem:rho-scaling}
Let \(\rho\in\F^\times\), and let \(\F'/\F\) be a field extension containing
an element \(s\) such that
\[
        s^2\rho_0=\rho .
\]
There is an algebra isomorphism
\[
\Theta_s:
\Aq(\rho)\otimes_\F\F'
\longrightarrow
\Aq(\rho_0)\otimes_\F\F'
\]
determined by
\begin{align*}
\cW_{-n}^{(\rho)}&\longmapsto s\,\cW_{-n}^{(\rho_0)},
&
\cW_{n+1}^{(\rho)}&\longmapsto s\,\cW_{n+1}^{(\rho_0)},\\
\cG_{n+1}^{(\rho)}&\longmapsto s^2\cG_{n+1}^{(\rho_0)},
&
\ctG_{n+1}^{(\rho)}&\longmapsto s^2\ctG_{n+1}^{(\rho_0)}.
\end{align*}
Moreover, for every \(n\geq1\),
\begin{equation}
\label{eq:Delta-scaling}
        \Theta_s(\Delta_n^{(\rho)})
        =
        s^2\Delta_n^{(\rho_0)} .
\end{equation}
\end{lemma}

\begin{proof}
The defining relations of the current algebra are homogeneous with respect
to the following weights:
\[
        \deg_s\cW_{\pm}=1,
        \qquad
        \deg_s\cG=\deg_s\ctG=2,
        \qquad
        \deg_s\rho=2.
\]
Relations involving two \(W\)-generators have total \(s\)-degree \(2\);
relations involving one \(W\)-generator and one \(G\)-type generator have
total \(s\)-degree \(3\); and relations involving two \(G\)-type generators
have total \(s\)-degree \(4\).  The only parameter-dependent terms have the
form \(\rho\,\cW\).  Under the proposed assignment, such a term is sent to
\[
        \rho\,s\,\cW^{(\rho_0)}
        =
        s^3\rho_0\,\cW^{(\rho_0)},
\]
which is exactly the scaling of the corresponding relation in
\(\Aq(\rho_0)\).  Hence all defining relations are preserved.  Applying the
same construction with \(s^{-1}\) gives the inverse isomorphism.

It remains to check the central elements.  By the quantum-determinant
formula of \cite[Proposition~2.1]{BaseilhacBelliard2017}, and equivalently
by its coefficient form in \cite[Lemma~2.1]{BaseilhacBelliard2017}, each
coefficient \(\Delta_n^{(\rho)}\) is a sum of terms of the following
homogeneous types:
\[
        \cG+\ctG,\qquad
        \cW\cW,\qquad
        \rho^{-1}\cG\ctG .
\]
Under \(\Theta_s\), the first two types are multiplied by \(s^2\).  The last
type is also multiplied by \(s^2\), because
\[
        \rho^{-1}s^4
        =
        (s^2\rho_0)^{-1}s^4
        =
        s^2\rho_0^{-1}.
\]
Therefore every term in \(\Delta_n^{(\rho)}\) is sent to \(s^2\) times the
corresponding term in \(\Delta_n^{(\rho_0)}\), proving
\eqref{eq:Delta-scaling}.
\end{proof}

\begin{theorem}[Baseilhac--Belliard basis, arbitrary parameter]
\label{thm:BB-general}
Assume that \(\operatorname{char}\F\neq2\).  Let
\(\rho\in\F^\times\), and let
\(\boldsymbol\delta=(\delta_1,\delta_2,\ldots)\) be any scalar sequence.
Then the monomials of the form \eqref{eq:BB-monomials}, interpreted in the
quotient \(\BBqquot{\rho}{\boldsymbol\delta}\), form an \(\F\)-basis of
\(\BBqquot{\rho}{\boldsymbol\delta}\).  Consequently,
\cite[Conjecture~1]{BaseilhacBelliard2017} holds.
\end{theorem}

\begin{proof}
By Lemma~\ref{lem:nonvanishing-scalars}, $\rho_0\neq0$.  Hence there is
an algebraic field extension $\F'/\F$ containing an element $s$ with
$s^2\rho_0=\rho$.  By Lemma~\ref{lem:rho-scaling}, after extending
scalars to \(\F'\), the quotient
\(\BBqquot{\rho}{\boldsymbol\delta}\) is isomorphic to the normalised
quotient
\[
        \BBqquot{\rho_0}{\boldsymbol\delta/s^2},
        \qquad
        \boldsymbol\delta/s^2
        =
        (\delta_1/s^2,\delta_2/s^2,\ldots).
\]
Indeed, the relation
\(\Delta_n^{(\rho)}=2\delta_n\) is sent to
\[
        s^2\Delta_n^{(\rho_0)}=2\delta_n,
\]
or equivalently
\[
        \Delta_n^{(\rho_0)}=2\delta_n/s^2.
\]

By Theorem~\ref{thm:BB-normalised}, the monomials
\eqref{eq:BB-monomials} form a basis of the normalised quotient with central
values \(\boldsymbol\delta/s^2\).  Under \(\Theta_s\), a monomial
\eqref{eq:BB-monomials} is multiplied by a non-zero power of \(s\), namely
\[
        s^{a_1+\cdots+a_N+c_1+\cdots+c_M
          +2(b_1+\cdots+b_P)}.
\]
Therefore the corresponding monomials form a basis of
\[
        \BBqquot{\rho}{\boldsymbol\delta}\otimes_\F\F' .
\]

It remains to descend the basis property from \(\F'\) to \(\F\).  Linear
independence descends immediately: any non-trivial finite relation over
\(\F\) would remain a non-trivial relation after tensoring with \(\F'\).
For spanning, let \(V\) be the \(\F\)-span of the monomials
\eqref{eq:BB-monomials} inside \(\BBqquot{\rho}{\boldsymbol\delta}\).  We
have shown that
\[
        ( \BBqquot{\rho}{\boldsymbol\delta}/V )
        \otimes_\F \F' =0 .
\]
Since \(\F'/\F\) is a faithfully flat field extension, this implies
\(\BBqquot{\rho}{\boldsymbol\delta}/V=0\).  Hence the monomials span over
\(\F\), and therefore form an \(\F\)-basis.
\end{proof}

\begin{remark}
The argument proves more than the single order stated in
\cite[Conjecture~1]{BaseilhacBelliard2017}.  Every quotient
\(\BBqquot{\rho}{\boldsymbol\delta}\), with
\(\rho\neq0\) and \(\operatorname{char}\F\neq2\), inherits all twelve PBW
bases obtained from the two triples
\[
        (W^-,G,W^+),
        \qquad
        (W^-,\widetilde G,W^+)
\]
and the six block orders of Theorem~\ref{thm:central-fibre-PBW}.
\end{remark}





\section{A large-index triangular crossing theorem}
\label{sec:large-crossing}

Let $U$ denote the image of $\Uplus$ in the $q$-shuffle algebra.  In this
section only, we omit the superscript $\sh$ from the alternating words and
write
\[
 W_{-n}=(xy)^nx,\qquad
 W_{n+1}=(yx)^ny,\qquad
 G_n=(yx)^n,\qquad
 \widetilde G_n=(xy)^n.
\]
All products in this section are $q$-shuffle products, denoted by $\star$.

We use the alternating PBW order
\[
        W^+<W^-<\widetilde G,
\]
with increasing index order inside each block.  For finite non-decreasing
sequences
\[
 \lambda=(\lambda_1,\ldots,\lambda_r),\quad \lambda_i\geq1,
 \qquad
 \mu=(\mu_1,\ldots,\mu_s),\quad \mu_i\geq0,
\]
\[
 \nu=(\nu_1,\ldots,\nu_t),\quad \nu_i\geq1,
\]
put
\[
 P_\lambda=W_{\lambda_1}\star\cdots\star W_{\lambda_r},\qquad
 Q_\mu=W_{-\mu_1}\star\cdots\star W_{-\mu_s},
\]
\[
 H_\nu=\widetilde G_{\nu_1}\star\cdots\star
       \widetilde G_{\nu_t}.
\]
The products $P_\lambda\star Q_\mu\star H_\nu$ form a basis of $U$ by
Theorem~\ref{thm:qshuffle-twelve}.

For a finite multiset $A$ of positive integers, define
\[
 \operatorname{prof}(A)
 =\bigl(|A|;a_{|A|},a_{|A|-1},\ldots,a_1\bigr),
\]
where $a_1\leq\cdots\leq a_{|A|}$ are its entries, and order profiles
lexicographically.  We use this order for both positive and middle blocks.

\begin{lemma}[Profile maps]
\label{lem:profile-maps}
Let $A,B,C$ be finite multisets of positive integers.
\begin{enumerate}[label=\textup{(\roman*)}]
\item If $\operatorname{prof}(A)<\operatorname{prof}(B)$, then
\[
 \operatorname{prof}(A\sqcup C)<\operatorname{prof}(B\sqcup C).
\]
\item Suppose $A$ and $B$ are non-empty and $N$ is larger than every entry
of $A$ and $B$.  Define
\[
 T_N(A)=(A\setminus\{\max A\})\sqcup\{N+\max A\}.
\]
Then $T_N$ is injective and strictly order-preserving for the profile order.
\end{enumerate}
\end{lemma}

\begin{proof}
For (i), compare multiplicities from the largest integer downwards; adding
the same multiset changes both multiplicity functions by the same amount.
For (ii), $N+\max A$ is the unique largest entry of $T_N(A)$, so $\max A$
and then $A$ are recovered from the image.  If $\max A<\max B$, the unique
largest entries of the two images preserve this inequality.  If the maxima
are equal, remove one copy of the common maximum from both multisets; after
the common new largest entry is ignored, the original comparison of the
remaining entries is unchanged.
\end{proof}

\begin{lemma}[Large-index leading crossings]
\label{lem:large-crossings}
The following statements hold.
\begin{enumerate}[label=\textup{(\roman*)}]
\item For $m\geq1$,
\begin{equation}
\label{eq:G-leading-positive}
 G_m=W_m\star W_0+
 \{\text{terms of lower positive profile}\}.
\end{equation}
\item For $m\geq1$,
\begin{equation}
\label{eq:W0-Wplus-leading}
 W_0\star W_m=q^{-2}W_m\star W_0+
 \{\text{terms of lower positive profile}\}.
\end{equation}
\item If $a\geq1$ and $N\geq a-1$, then
\begin{equation}
\label{eq:large-W-crossing}
 [W_a,W_{-N}]
 =(1-q^{-2})W_{N+a}\star W_0+
 \{\text{terms of lower positive profile}\}.
\end{equation}
\item If $i\geq1$ and $N\geq i$, then
\begin{equation}
\label{eq:large-tG-crossing}
 [\widetilde G_i,W_{-N}]
 =(q^{-2}-1)W_0\star\widetilde G_{N+i}+
 \{\text{terms of lower middle profile}\}.
\end{equation}
\item For $i\geq1$,
\begin{equation}
\label{eq:tG-W0-exact}
 \widetilde G_i\star W_0
 =q^{-2}W_0\star\widetilde G_i+(1-q^{-2})W_{-i}.
\end{equation}
\end{enumerate}
\end{lemma}

\begin{proof}
Let
\[
 \widetilde G(t)=\sum_{n\geq0}\widetilde G_nt^n,
 \qquad D(t)=\sum_{n\geq0}D_nt^n,
\]
where $D_0=1$ and
$D(t)\star\widetilde G(t)=1=\widetilde G(t)\star D(t)$; see
\cite[Definition~9.5, Definition~9.11, and Lemma~9.12]
{Terwilliger2019Alternating}.  Each $D_j$ is a polynomial in
$\widetilde G_1,\ldots,\widetilde G_j$.  By
\cite[Theorem~9.15, equation~(70)]{Terwilliger2019Alternating},
\begin{equation}
\label{eq:G-D-expansion}
 G_m=q^{-2m}D_m+
 \sum_{a+b+c+1=m}W_{c+1}\star D_b\star W_{-a}.
\end{equation}
The unique summand with largest positive index is obtained from
$(a,b,c)=(0,0,m-1)$ and equals $W_m\star W_0$.  Reordering $D_b$ past a
negative word creates no positive $W$-factor, by the crossing formulas in
\cite[Lemma~15.1]{Terwilliger2019Alternating}.  This proves (i).

Proposition~5.7 of the same reference gives
\[
 [W_0,W_m]=(1-q^{-2})(\widetilde G_m-G_m).
\]
Together with (i), this gives (ii).

Write $a=i+1$.  For $i\leq N$, Lemma~16.1 of the same reference gives
\begin{equation}
\label{eq:exact-W-large-crossing}
 \begin{split}
 [W_{i+1},W_{-N}]
 ={}&(1-q^{-2})\sum_{\ell=0}^{i}
       G_{N+1+\ell}\star\widetilde G_{i-\ell}\\
 &-(1-q^{-2})\sum_{\ell=0}^{i}
       G_{i-\ell}\star\widetilde G_{N+1+\ell}.
 \end{split}
\end{equation}
By (i), the unique largest positive index is $N+i+1=N+a$; it occurs in
the first sum at $\ell=i$.  This proves (iii).

For $i\leq N$, Lemma~15.1 gives
\begin{equation}
\label{eq:exact-tG-large-crossing}
 \begin{split}
 \widetilde G_i\star W_{-N}
 ={}&W_{-N}\star\widetilde G_i
 +(1-q^{-2})\sum_{\ell=1}^{i}
   W_{-N-\ell}\star\widetilde G_{i-\ell}\\
 &+(q^{-2}-1)\sum_{\ell=1}^{i}
   W_{\ell-i}\star\widetilde G_{N+\ell}.
 \end{split}
\end{equation}
The unique term with middle index $N+i$ is the term $\ell=i$ in the final
sum, proving (iv).  Taking the case $N=0<i$ in the complementary part of
Lemma~15.1 gives (v).
\end{proof}

\begin{lemma}[Leading crossings of blocks]
\label{lem:block-leading-crossings}
Let $N$ be larger than every index in the relevant block.
\begin{enumerate}[label=\textup{(\roman*)}]
\item Let $\lambda$ be non-empty, let $p=\max\lambda$, and let $e$ be the
multiplicity of $p$.  Then
\begin{equation}
\label{eq:positive-block-crossing}
 [P_\lambda,W_{-N}]
 =(1-q^{-2e})P_{T_N(\lambda)}\star W_0
 +\{\text{terms of lower positive profile}\}.
\end{equation}
\item Let $\nu$ be non-empty, let $i=\max\nu$, let $e$ be its multiplicity,
and put $t=|\nu|$.  Then
\begin{equation}
\label{eq:middle-block-crossing}
 [H_\nu,W_{-N}]
 =q^{-2(t-e)}(q^{-2e}-1)
 W_0\star H_{T_N(\nu)}
 +\{\text{terms of lower middle profile}\}.
\end{equation}
\end{enumerate}
\end{lemma}

\begin{proof}
Expand each commutator by the Leibniz rule.  For (i), crossing a factor
$W_a$ with $W_{-N}$ can create a positive index $N+a$ by
\eqref{eq:large-W-crossing}.  The largest target is therefore obtained only
from the $e$ occurrences of $W_p$.  If the selected occurrence has $s$
further copies of $W_p$ to its right, the newly created $W_0$ must cross
these copies.  Equation~\eqref{eq:W0-Wplus-leading} contributes
$q^{-2s}$ to the leading profile; the positive words commute among
themselves.  Hence the leading coefficient is
\[
 (1-q^{-2})\sum_{s=0}^{e-1}q^{-2s}=1-q^{-2e}.
\]
All other choices and all correction terms have lower positive profile.

For (ii), only the $e$ occurrences of the maximal middle index $i$ can
create $N+i$.  The numbers of middle factors to the left of these
occurrences are $t-e,t-e+1,\ldots,t-1$.  Moving the new $W_0$ to the left
by \eqref{eq:tG-W0-exact} contributes the leading factors
$q^{-2s}$; every correction in that move decreases the middle length.  The
leading coefficient is therefore
\[
 (q^{-2}-1)\sum_{s=t-e}^{t-1}q^{-2s}
 =q^{-2(t-e)}(q^{-2e}-1).
\]
This proves both assertions.
\end{proof}

\begin{theorem}[Negative-tail centraliser]
\label{thm:tail-centraliser}
For every $N_0\in\N$,
\begin{equation}
\label{eq:tail-centraliser}
 \Cent_U\{W_{-N}\mid N\geq N_0\}
 =\F[W_0,W_{-1},W_{-2},\ldots].
\end{equation}
\end{theorem}

\begin{proof}
The inclusion from right to left follows from the pairwise commutativity of
the negative alternating words.  Put
\[
        \mathcal N=\F[W_0,W_{-1},W_{-2},\ldots].
\]
Theorem~\ref{thm:qshuffle-twelve} and the commutativity of the negative
family show that $\mathcal N$ is a polynomial algebra, hence a domain.

Let $z$ centralise $W_{-N}$ for every $N\geq N_0$.  Its PBW expansion can
be grouped uniquely as
\begin{equation}
\label{eq:grouped-tail-expansion}
        z=\sum_{\lambda,\nu}P_\lambda\star
        f_{\lambda,\nu}\star H_\nu,
        \qquad f_{\lambda,\nu}\in\mathcal N,
\end{equation}
with only finitely many non-zero coefficients.

Suppose that a non-empty positive block occurs.  Choose a maximal positive
profile among the non-zero terms, represented by $\lambda$, and choose
$N\geq N_0$ larger than every positive and middle index occurring in
\eqref{eq:grouped-tail-expansion}.  Let $p=\max\lambda$ and let $e$ be its
multiplicity.  By Lemma~\ref{lem:block-leading-crossings}, the component of
$[z,W_{-N}]$ with target positive profile $T_N(\lambda)$ is
\[
 (1-q^{-2e})P_{T_N(\lambda)}\star
 W_0f_{\lambda,\nu}\star H_\nu
\]
for each middle block $\nu$.  A commutator originating in a middle block
keeps the old positive profile.  A smaller source positive profile cannot
produce $T_N(\lambda)$, by the strict order preservation in
Lemma~\ref{lem:profile-maps}; injectivity also excludes a different source
with the same target.  Thus this component cannot cancel.  Its coefficient
is non-zero because $1-q^{-2e}\neq0$ and multiplication by $W_0$ is
injective in the domain $\mathcal N$.  This contradicts
$[z,W_{-N}]=0$.  Hence no positive block occurs.

We may now write $z=\sum_\nu f_\nu\star H_\nu$.  If a non-empty middle
block occurs, choose one with maximal middle profile, and choose $N\geq N_0$
larger than every middle index occurring in $z$.  If its maximal index has
multiplicity $e$ and its length is $t$, equation
\eqref{eq:middle-block-crossing} produces the unique maximal target profile
with coefficient
\[
        q^{-2(t-e)}(q^{-2e}-1)W_0f_\nu,
\]
which is non-zero.  Lemma~\ref{lem:profile-maps} again excludes
cancellation from another source.  This contradicts centrality.  Therefore
no middle block occurs, and $z\in\mathcal N$.
\end{proof}

\begin{corollary}
\label{cor:conj167}
One has
\[
 \Cent_{\Oq}\bigl(\F[W_0,W_{-1},W_{-2},\ldots]\bigr)
 =\F[W_0,W_{-1},W_{-2},\ldots].
\]
In particular, Terwilliger's Conjecture~16.7 holds.
\end{corollary}

\begin{proof}
Let $H=\F[W_0,W_{-1},W_{-2},\ldots]\subseteq\Oq$.  By
Theorem~\ref{thm:twelve-PBW}, $H$ is a polynomial algebra.  Moreover,
the negative-family monomials have distinct, non-zero initial forms by
Theorem~\ref{thm:all-leading}; hence the induced filtration on $H$ has
\[
 \gr H=\F[W_0^{\sh},W_{-1}^{\sh},W_{-2}^{\sh},\ldots]\subseteq U.
\]
Theorem~\ref{thm:tail-centraliser}, with $N_0=0$, shows that this graded
subalgebra is self-centralising in $U$.  Lemma~\ref{lem:filtered-centraliser}
then yields the displayed equality in $\Oq$.
\end{proof}


\section{Centralisers of alternating subalgebras}
\label{sec:centralisers}

\subsection{The imaginary-root centraliser}

For $r\in\Z$, define a two-sided real-root chain by
\[
 R_r=
 \begin{cases}
 B_{r\delta+\alpha_1}, & r\geq0,\\[1mm]
 B_{(-r-1)\delta+\alpha_0}, & r\leq-1,
 \end{cases}
 \qquad
 I_n=B_{n\delta}\quad(n\geq1).
\]
Thus $R_0=W_1$ and $R_{-1}=W_0$.  Order the root vectors by
\begin{equation}
\label{eq:BK-centraliser-order}
 R_r<R_s\quad\Longleftrightarrow\quad r>s,
 \qquad R_r<I_m,
 \qquad I_m<I_n\quad\Longleftrightarrow\quad m<n.
\end{equation}
By Theorem~\ref{thm:BK-nonroot}, this total order gives a PBW basis of
$\Oq$.  Consequently every element has a unique expansion
\begin{equation}
\label{eq:BK-centraliser-expansion}
 \sum_{\boldsymbol r}
 R_{r_1}\cdots R_{r_d}
 f_{\boldsymbol r}(I_1,I_2,\ldots),
 \qquad r_1\geq r_2\geq\cdots\geq r_d.
\end{equation}

Put
\[
        \kappa=q^{-1}(q-q^{-1})^2[2]_q.
\]
In the normalisation of this paper, equations (3.9)--(3.10), (5.8) with $m=1$,
and Corollary~5.11 of \cite{BaseilhacKolb2020} give the following
relations.  Their proofs are formal consequences of the defining relations;
after substituting the present value of $c$, the displayed identities have
Laurent-polynomial coefficients in $q$, so the same calculations apply over
our field at the fixed non-root parameter:
\begin{align}
[I_1,R_r]&=\kappa(R_{r-1}-R_{r+1}),
\label{eq:I1-real-shift}\\
R_rR_{r+1}&=q^{-2}R_{r+1}R_r-I_1,
\label{eq:adjacent-real-straightening}\\
[I_m,I_n]&=0.
\label{eq:imaginary-commute}
\end{align}
Indeed, the parameter $c$ in that reference is
$c=q^{-1}(q-q^{-1})^2$ for the Dolan--Grady normalisation
\eqref{eq:qDG0}--\eqref{eq:qDG1}.

For a non-empty non-increasing integer sequence
$\boldsymbol r=(r_1,\ldots,r_d)$, define
\[
        T(\boldsymbol r)
        =\operatorname{sort}_{\geq}(r_1+1,r_2,\ldots,r_d).
\]
Lexicographic order is taken after the length $d$.  The map $T$ is injective
and strictly order-preserving: its largest entry is one more than the
largest entry of the source, and after deleting these entries the remaining
comparison is unchanged.

\begin{lemma}[Leading real profile]
\label{lem:leading-real-profile}
Let $r_1\geq\cdots\geq r_d$, let $r=r_1$, and suppose that $r$ has
multiplicity $e$.  Then
\begin{equation}
\label{eq:leading-real-commutator}
 \begin{split}
 [I_1,R_{r_1}\cdots R_{r_d}]
 ={}&-\kappa\left(\sum_{j=0}^{e-1}q^{-2j}\right)
 R_{r+1}R_r^{e-1}R_{r_{e+1}}\cdots R_{r_d}\\
 &+\{\text{terms of lower real profile or with fewer real factors}\}.
 \end{split}
\end{equation}
\end{lemma}

\begin{proof}
Use the Leibniz rule and \eqref{eq:I1-real-shift}.  Only the raising terms
$-\kappa R_{r+1}$ coming from the first $e$ factors can produce the profile
$T(\boldsymbol r)$.  If the selected occurrence has $j$ copies of $R_r$ to
its left, relation \eqref{eq:adjacent-real-straightening} contributes
$q^{-2j}$ while moving $R_{r+1}$ to the left.  Every correction $-I_1$ in
that move removes two real factors.  Raising a smaller entry or using a
lowering term gives a lower profile.  Summing over $j=0,\ldots,e-1$ proves
the formula.
\end{proof}

\begin{lemma}[Imaginary-root centraliser]
\label{lem:imaginary-centraliser}
One has
\begin{equation}
\label{eq:I1-centraliser}
        \Cent_{\Oq}(I_1)=\F[I_1,I_2,I_3,\ldots].
\end{equation}
\end{lemma}

\begin{proof}
The inclusion from right to left follows from
\eqref{eq:imaginary-commute}.  Let $x\in\Cent_{\Oq}(I_1)$ and use the PBW
expansion \eqref{eq:BK-centraliser-expansion}.  Suppose that a real factor
occurs.  Among all terms with non-zero coefficient, choose one with maximal
real profile $(d;r_1,\ldots,r_d)$, and let the maximal entry $r=r_1$ have
multiplicity $e$.

By Lemma~\ref{lem:leading-real-profile}, its contribution to the profile
$T(r_1,\ldots,r_d)$ in $[I_1,x]$ has coefficient
\[
 -\kappa\sum_{j=0}^{e-1}q^{-2j}
 =-\kappa\frac{1-q^{-2e}}{1-q^{-2}},
\]
which is non-zero.  A smaller source profile cannot give the same target,
because $T$ is strictly order-preserving and injective.  A lowering term
cannot create the new largest entry $r+1$, and terms with fewer real factors
remain lower in the length-first order.  The non-zero polynomial
$f_{\boldsymbol r}(I_1,I_2,\ldots)$ is unchanged in this leading component.
Thus the component cannot cancel, contradicting $[I_1,x]=0$.  Hence no real
factor occurs, and $x$ belongs to the displayed polynomial algebra.
\end{proof}

\begin{corollary}
\label{cor:conj168}
One has
\begin{equation}
\label{eq:tG1-centraliser}
 \Cent_{\Oq}(\widetilde G_1)
 =\F[\widetilde G_1,\widetilde G_2,\widetilde G_3,\ldots].
\end{equation}
In particular, Terwilliger's Conjecture~16.8 holds.
\end{corollary}

\begin{proof}
In the recursion \eqref{eq:root-tG-recursion}, the coefficient of
$\widetilde G_n$ is $[2n]_q[2]_q^nB_{0\delta}$, while the term containing
$I_n=B_{n\delta}$ linearly is $[n]_qI_n\widetilde G_0$.  Both coefficients
are non-zero.  Induction on $n$, using the commutativity of the $I_m$, shows
that the two families are related by a triangular polynomial change of
variables with non-zero diagonal coefficients.  Therefore
\[
 \F[I_1,I_2,\ldots]
 =\F[\widetilde G_1,\widetilde G_2,\ldots].
\]
Equation~\eqref{eq:tG1-Bdelta} gives $\widetilde G_1=-qI_1$, so the two
first elements have the same centraliser.  The result follows from
Lemma~\ref{lem:imaginary-centraliser}.
\end{proof}

\begin{corollary}
\label{cor:four-maximal}
The following four polynomial subalgebras of $\Oq$ are self-centralising,
and hence maximal commutative:
\[
 \F[W_0,W_{-1},W_{-2},\ldots],
 \qquad
 \F[W_1,W_2,W_3,\ldots],
\]
\[
 \F[G_1,G_2,G_3,\ldots],
 \qquad
 \F[\widetilde G_1,\widetilde G_2,\widetilde G_3,\ldots].
\]
Moreover,
\[
 \Cent_{\Oq}(\widetilde G_1)
 =\F[\widetilde G_1,\widetilde G_2,\ldots],
 \qquad
 \Cent_{\Oq}(G_1)=\F[G_1,G_2,\ldots].
\]
\end{corollary}

\begin{proof}
The negative algebra is self-centralising by
Corollary~\ref{cor:conj167}; applying $\sigma$ gives the positive algebra.
Let
$\widetilde{\mathcal H}=\F[\widetilde G_1,\widetilde G_2,\ldots]$.
Since this algebra is commutative and contains $\widetilde G_1$,
Corollary~\ref{cor:conj168} gives
\[
 \widetilde{\mathcal H}
 \subseteq\Cent_{\Oq}(\widetilde{\mathcal H})
 \subseteq\Cent_{\Oq}(\widetilde G_1)
 =\widetilde{\mathcal H}.
\]
Thus it is self-centralising.  Applying $\dagger$ gives the corresponding
claims for the $G$-family, including the formula for $\Cent_{\Oq}(G_1)$.
A self-centralising commutative subalgebra is maximal commutative.
\end{proof}


\begin{thebibliography}{99}

\bibitem{AppelVlaar2022}
A.~Appel and B.~Vlaar,
\emph{Universal $K$-matrices for quantum Kac--Moody algebras},
Represent. Theory \textbf{26} (2022), 764--824.

\bibitem{Baseilhac2005}
P.~Baseilhac,
\emph{Deformed Dolan--Grady relations in quantum integrable models},
Nuclear Phys. B \textbf{709} (2005), 491--521.

\bibitem{BaseilhacBelliard2010}
P.~Baseilhac and S.~Belliard,
\emph{Generalized $q$-Onsager algebras and boundary affine Toda field theories},
Lett. Math. Phys. \textbf{93} (2010), 213--228.

\bibitem{BaseilhacBelliard2013}
P.~Baseilhac and S.~Belliard,
\emph{The half-infinite $XXZ$ chain in Onsager's approach},
Nuclear Phys. B \textbf{873} (2013), 550--584.

\bibitem{BaseilhacBelliard2017}
P.~Baseilhac and S.~Belliard,
\emph{An attractive basis for the $q$-Onsager algebra},
arXiv:1704.02950.

\bibitem{BaseilhacGainutdinovLemarthe2026}
P.~Baseilhac, A.~M.~Gainutdinov and G.~Lemarthe,
\emph{Universal TT- and TQ-relations via centrally extended $q$-Onsager algebra},
Nuclear Phys. B \textbf{1026} (2026), 117427.

\bibitem{BaseilhacKoizumi2005}
P.~Baseilhac and K.~Koizumi,
\emph{A new (in)finite-dimensional algebra for quantum integrable models},
Nuclear Phys. B \textbf{720} (2005), 325--347.

\bibitem{BaseilhacKolb2020}
P.~Baseilhac and S.~Kolb,
\emph{Braid group action and root vectors for the $q$-Onsager algebra},
Transform. Groups \textbf{25} (2020), 363--389.

\bibitem{BaseilhacShigechi2010}
P.~Baseilhac and K.~Shigechi,
\emph{A new current algebra and the reflection equation},
Lett. Math. Phys. \textbf{92} (2010), 47--65.

\bibitem{Beck1994}
J.~Beck,
\emph{Convex bases of PBW type for quantum affine algebras},
Comm. Math. Phys. \textbf{165} (1994), 193--199.

\bibitem{Damiani1993}
I.~Damiani,
\emph{A basis of type Poincar\'e--Birkhoff--Witt for the quantum algebra of
$\widehat{\mathfrak{sl}}_2$},
J. Algebra \textbf{161} (1993), 291--310.

\bibitem{Damiani1998}
I.~Damiani,
\emph{La $R$-matrice pour les alg\`ebres quantiques de type affine non tordu},
Ann. Sci. \`Ecole Norm. Sup. (4) \textbf{31} (1998), 493--523.

\bibitem{DingFrenkel1993}
J.~Ding and I.~B.~Frenkel,
\emph{Isomorphism of two realizations of quantum affine algebra
$U_q(\widehat{\mathfrak{gl}}(n))$},
Comm. Math. Phys. \textbf{156} (1993), 277--300.

\bibitem{DolanGrady1982}
L.~Dolan and M.~Grady,
\emph{Conserved charges from self-duality},
Phys. Rev. D \textbf{25} (1982), 1587--1604.

\bibitem{KhoroshkinTolstoy1991}
S.~M.~Khoroshkin and V.~N.~Tolstoy,
\emph{Universal $R$-matrix for quantized (super)algebras},
Comm. Math. Phys. \textbf{141} (1991), 599--617.

\bibitem{Kolb2014}
S.~Kolb,
\emph{Quantum symmetric Kac--Moody pairs},
Adv. Math. \textbf{267} (2014), 395--469.

\bibitem{LemartheBaseilhacGainutdinov2026}
G.~Lemarthe, P.~Baseilhac and A.~M.~Gainutdinov,
\emph{Fused $K$-operators and the $q$-Onsager algebra},
SIGMA Symmetry Integrability Geom. Methods Appl. \textbf{22} (2026), 026,
74 pp.

\bibitem{Letzter2002}
G.~Letzter,
\emph{Coideal subalgebras and quantum symmetric pairs},
in \emph{New Directions in Hopf Algebras}, Math. Sci. Res. Inst. Publ.
\textbf{43}, Cambridge Univ. Press, Cambridge, 2002, 117--165.

\bibitem{LuRuanWang2023}
M.~Lu, S.~Ruan and W.~Wang,
\emph{$\imath$Hall algebra of the projective line and $q$-Onsager algebra},
Trans. Amer. Math. Soc. \textbf{376} (2023), 1475--1505.

\bibitem{LuWang2021}
M.~Lu and W.~Wang,
\emph{A Drinfeld type presentation of affine $\imath$quantum groups I:
split ADE type},
Adv. Math. \textbf{393} (2021), 108111, 46 pp.

\bibitem{Onsager1944}
L.~Onsager,
\emph{Crystal statistics. I. A two-dimensional model with an order-disorder transition},
Phys. Rev. \textbf{65} (1944), 117--149.

\bibitem{Rosso1998}
M.~Rosso,
\emph{Quantum groups and quantum shuffles},
Invent. Math. \textbf{133} (1998), 399--416.

\bibitem{Sklyanin1988}
E.~K.~Sklyanin,
\emph{Boundary conditions for integrable quantum systems},
J. Phys. A \textbf{21} (1988), 2375--2389.

\bibitem{Terwilliger2001}
P.~Terwilliger,
\emph{Two relations that generalize the $q$-Serre relations and the Dolan--Grady relations},
in \emph{Physics and Combinatorics 1999}, World Scientific, River Edge, NJ,
2001, 377--398.

\bibitem{Terwilliger2017PositivePart}
P.~Terwilliger,
\emph{The $q$-Onsager algebra and the positive part of
$U_q(\widehat{\mathfrak{sl}}_2)$},
Linear Algebra Appl. \textbf{521} (2017), 19--56.

\bibitem{Terwilliger2019Catalan}
P.~Terwilliger,
\emph{Using Catalan words and a $q$-shuffle algebra to describe a PBW basis
for the positive part of $U_q(\widehat{\mathfrak{sl}}_2)$},
J. Algebra \textbf{525} (2019), 359--373.

\bibitem{Terwilliger2019Alternating}
P.~Terwilliger,
\emph{The alternating PBW basis for the positive part of
$U_q(\widehat{\mathfrak{sl}}_2)$},
J. Math. Phys. \textbf{60} (2019), 071704, 27 pp.

\bibitem{Terwilliger2021ACE}
P.~Terwilliger,
\emph{The alternating central extension of the $q$-Onsager algebra},
Comm. Math. Phys. \textbf{387} (2021), 1771--1819.

\bibitem{Terwilliger2022OqACE}
P.~Terwilliger,
\emph{The $q$-Onsager algebra and its alternating central extension},
Nuclear Phys. B \textbf{975} (2022), 115662, 37 pp.


\end{thebibliography}

\end{document}
